\chapter{常微分方程与动力系统初步}
\label{ch:ode}

% ════════════════════════════════════════════════════════════════
%  数学基础部：为 SLAM/估计/规划/控制各部打动力系统地基。
%  主线 synth 自 Tutorial 120_常微分方程；联网核 矩阵指数闭式/稳定性判据。
%  自包含独立记录，遵循 v5.0 理论教学规范（动机→反面→历史→理论→陷阱→练习
%  + 符号表/速查/后续关系）。TikZ 内联重画，记号与 lie_theory/rigid_body 一致。
% ════════════════════════════════════════════════════════════════

\begin{practice}[前置自测]
开始之前，先试答下列问题。若已能流畅作答，可只速读本章表格与速查卡；若卡住，正是本章要为你解决的。
\begin{enumerate}
  \item 一根受力矩驱动的单连杆机械臂满足 $I\ddot\theta+b\dot\theta+mg\ell\sin\theta=\tau$。给定初始角 $\theta(0)$ 与角速度 $\dot\theta(0)$，它的运动 $\theta(t)$ 是否被\emph{唯一}确定？凭什么？
  \item 把上面这个二阶方程改写成一阶向量形式 $\dot{\mathbf{x}}=\mathbf{f}(\mathbf{x},\tau)$，状态 $\mathbf{x}$ 应该取什么？只取 $\theta$ 行不行？
  \item 常系数线性系统 $\dot{\mathbf{x}}=\mathbf{A}\mathbf{x}$ 的解为什么是 $\mathbf{x}(t)=e^{t\mathbf{A}}\mathbf{x}_0$？矩阵指数 $e^{t\mathbf{A}}$ 到底是什么？
  \item 你给机械臂写了 PD$+$重力补偿，想\emph{证明}它一定收敛到目标——但闭式解算不出来。除了仿真碰运气，还有别的办法吗？
  \item 对称正定矩阵 $\mathbf{P}\succ0$ 有哪些等价刻画？二次型 $V(\mathbf{x})=\mathbf{x}^\top\mathbf{P}\mathbf{x}$ 的梯度 $\nabla V$ 等于什么？
\end{enumerate}
\end{practice}

\section*{本章目标}
学完本章，你应当能够：
\begin{enumerate}
  \item \textbf{把}任意机器人动力学（欧拉--拉格朗日二阶系统）化为一阶状态空间 ODE $\dot{\mathbf{x}}=\mathbf{f}(\mathbf{x},\mathbf{u},t)$，并解释为什么“每一次仿真器 \texttt{step()} 都是在数值求解这个 ODE”；
  \item \textbf{独立复述} Picard--Lindelöf 存在唯一性定理的压缩映射证明，说明 Lipschitz 条件为何是唯一性的命脉、它在非光滑动力学（滑模、接触）中如何失效；
  \item \textbf{掌握}线性系统的基本解矩阵 $\boldsymbol{\Phi}$ 与常数变易公式，理解“自由响应 $+$ 卷积响应”的叠加结构；
  \item \textbf{熟练计算}矩阵指数 $e^{t\mathbf{A}}$（对角化 / Jordan 形 / Laplace），并推导罗德里格斯公式，把它与 \cref{ch:lie} 的 $\SO(3)$ 指数映射焊接；
  \item \textbf{读懂}相平面、不动点分类与 Hartman--Grobman 定理，理解“双曲平衡点处线性化定性可靠”为何是局部 LQR 的依据；
  \item \textbf{运用} Lyapunov 直接法与 LaSalle 不变原理证明机器人控制律的稳定性（含 PD$+$重力补偿的经典证明）；
  \item \textbf{建立} Hurwitz 矩阵与 Lyapunov 方程 $\mathbf{A}^\top\mathbf{P}+\mathbf{P}\mathbf{A}=-\mathbf{Q}$ 的等价性，理解它为何是 LQR Riccati 方程之根；
  \item \textbf{区分}稳定、渐近稳定、吸引子、极限环等定性概念，为后续控制部与滤波连续模型打前置直觉。
\end{enumerate}

\subsection*{本章知识导航}
本章是一条“\emph{机器人连续运动的数学是什么、它的解有哪些性质、如何分析}”的主线。结构如下：从“ODE 是什么”出发，先问\emph{局部存在唯一}，再求解\emph{唯一能完全解出的线性一类}（矩阵指数），最后转入\emph{非线性定性分析}（相图、稳定性）。

\begin{center}
\small
\begin{tikzpicture}[
  node distance=5mm and 7mm, >=Stealth,
  blk/.style={draw, rounded corners, align=center, font=\small, inner sep=3pt, fill=main!6, draw=main!60!black},
  grp/.style={draw, rounded corners, align=center, font=\small, inner sep=3pt, fill=second!6, draw=second!60!black},
  ln/.style={->, thick, gray!60!black}]
\node[blk] (ode) {ODE 与状态空间\\ $\dot{\mathbf{x}}=\mathbf{f}(t,\mathbf{x})$};
\node[blk, right=of ode] (exist) {存在唯一\\ Picard--Lindelöf};
\node[blk, right=of exist] (lin) {线性系统\\ $\boldsymbol{\Phi}$、常数变易};
\node[grp, below=8mm of ode] (expm) {矩阵指数\\ $e^{t\mathbf{A}}\to$ Rodrigues};
\node[grp, right=of expm] (phase) {相图 / 不动点\\ Hartman--Grobman};
\node[grp, right=of phase] (lyap) {Lyapunov\\ 直接法 $+$ LaSalle};
\node[blk, below=8mm of phase, fill=third!8, draw=third!60!black] (lyapeq) {Lyapunov 方程\\ $\mathbf{A}^\top\mathbf{P}+\mathbf{P}\mathbf{A}=-\mathbf{Q}$};
\node[blk, right=of lyapeq, fill=third!8, draw=third!60!black] (attr) {吸引子 / 极限环\\（定性图景）};
\draw[ln] (ode)--(exist); \draw[ln] (exist)--(lin);
\draw[ln] (lin.south) -- ++(0,-4mm) -| (expm.north);
\draw[ln] (expm)--(phase); \draw[ln] (phase)--(lyap);
\draw[ln] (lyap.south) -- (lyapeq.north);
\draw[ln] (lyapeq)--(attr);
\end{tikzpicture}
\end{center}

\noindent\textbf{推荐路径}：\cref{sec:ode-basic,sec:ode-picard} 是任何读者的主干（语言 $+$ 解的存在唯一）；\cref{sec:ode-linear,sec:ode-expm} 给出唯一能闭式求解的线性一类，其中矩阵指数（\cref{sec:ode-expm}）是 \cref{ch:lie} 指数映射的“线性原型”，强烈建议对照阅读；\cref{sec:ode-phase,sec:ode-lyapunov,sec:ode-lyapeq} 转入非线性定性分析与稳定性，是控制部 \cref{ch:lyapunov-basics}、\cref{ch:lyapunov-stability} 的前置直觉；\cref{sec:ode-attractor} 收束于吸引子的定性图景。

\subsection*{前置知识桥接}
本章把若干前置工具搬到\emph{动力系统}语境：
\begin{itemize}
  \item \textbf{完备度量空间与压缩映射}（Banach 不动点定理）——\cref{sec:ode-picard} 把“求解 ODE”重写为“求一个积分算子的不动点”，存在唯一性立刻从压缩映射原理掉出来。这是抽象不动点定理在\emph{函数空间}上最漂亮的一个应用。
  \item \textbf{线性代数}（\cref{ch:linalg}）——特征值、对角化、Jordan 标准形是 \cref{sec:ode-expm} 闭式计算 $e^{t\mathbf{A}}$ 的钥匙：幂零部分让无穷级数退化为有限多项式。
  \item \textbf{对称正定矩阵与二次型}（\cref{ch:linalg}）——Lyapunov 理论（\cref{sec:ode-lyapunov,sec:ode-lyapeq}）的核心对象是正定二次型 $V(\mathbf{x})=\mathbf{x}^\top\mathbf{P}\mathbf{x}$；$\mathbf{P}\succ0\iff$ 所有特征值正，且 $\nabla V=2\mathbf{P}\mathbf{x}$（$\mathbf{P}$ 对称时），本章反复调用。
  \item \textbf{刚体运动与旋转}（\cref{ch:rigid_body}）——\cref{sec:ode-expm} 的罗德里格斯公式与 \cref{ch:rigid_body} 给出的轴角$\to$旋转矩阵公式\emph{一模一样}，这里从“$e^{t\mathbf{A}}$ 解 $\dot{\mathbf{R}}=\mathbf{R}[\boldsymbol{\omega}]_\times$”的角度重新得到它。
\end{itemize}

\subsection*{如果跳过本章会怎样}
不学这一章，你会在四个具体场景中反复栽跟头：
\begin{itemize}
  \item \textbf{“仿真器给的轨迹我不敢信。”} 控制律含 $\operatorname{sign}(\cdot)$（滑模），破坏 Lipschitz 条件，从而破坏解的唯一性（\cref{sec:ode-picard}）——仿真器选哪条解取决于离散方式。不懂存在唯一性，无法区分“代码 bug”与“数学上本就病态”。
  \item \textbf{“姿态怎么积分？”} 给定角速度 $\boldsymbol{\omega}(t)$ 要把 $\mathbf{R}(t)\in\SO(3)$ 积出来。用欧拉角会遇万向锁；正确做法是矩阵指数（\cref{sec:ode-expm}），它把“在李代数里走一步”映成“流形上的真旋转”——这正是 IMU 预积分的内核。
  \item \textbf{“PD 增益我只能瞎调。”} 想证明 PD$+$重力补偿一定收敛，却只会仿真试错。其实有一个 40 年前就写好的两行 Lyapunov 证明（\cref{sec:ode-lyapunov}）告诉你它\emph{一定}渐近稳定。不懂 Lyapunov，稳定性对你就是玄学。
  \item \textbf{“线性系统稳不稳，除了算特征值还有别的判据吗？”} Lyapunov 方程（\cref{sec:ode-lyapeq}）用一个线性矩阵方程的可解性刻画 Hurwitz，解 $\mathbf{P}$ 还直接给出 LQR 的值函数。不懂它，你看不透“为什么 LQR 自带稳定性证明”。
\end{itemize}

\begin{table}[H]\centering\small
\caption{本章阅读方式建议}
\begin{tabular}{lll}
\toprule
阅读方式 & 时间 & 适合谁 \\
\midrule
精读（含全部证明与练习） & 6--8 小时 & 首次系统学 ODE 理论、要为控制/估计打地基 \\
速读（跳过证明细节） & 2--3 小时 & 有本科 ODE 基础、想建立机器人视角 \\
速查（只看表格与速查卡） & 20 分钟 & 已学过、回来查特定定理 / 公式 \\
\bottomrule
\end{tabular}
\end{table}

\begin{note}
\textbf{本章记号与约定.}\;
状态向量记 $\mathbf{x}(t)\in\mathbb{R}^n$（机器人中常为 $\mathbf{x}=(\mathbf{q},\dot{\mathbf{q}})$，即位置 $+$ 速度）；向量场（ODE 右端）记 $\mathbf{f}(t,\mathbf{x})$。
矩阵粗体大写 $\mathbf{A},\mathbf{P},\mathbf{Q}$；$\mathbf{A}^\top$ 转置；$\mathbf{P}\succ0$（$\mathbf{P}\succeq0$）表对称正定（半正定）。
矩阵指数记 $e^{t\mathbf{A}}=\exp(t\mathbf{A})$；$\boldsymbol{\omega}\in\mathbb{R}^3$ 的反对称矩阵记 $[\boldsymbol{\omega}]_\times\in\so(3)$（与 \cref{ch:rigid_body,ch:lie} 一致，满足 $[\boldsymbol{\omega}]_\times\mathbf{v}=\boldsymbol{\omega}\times\mathbf{v}$）。
机械臂动力学量沿用 \cref{ch:lie} 之前章节：惯性矩阵 $\mathbf{M}(\mathbf{q})\succ0$、科氏离心矩阵 $\mathbf{C}(\mathbf{q},\dot{\mathbf{q}})$、重力项 $\mathbf{G}(\mathbf{q})$、关节力矩 $\boldsymbol{\tau}$。
Lyapunov 函数记 $V(\mathbf{x})$，其沿轨迹导数记 $\dot V(\mathbf{x})=\nabla V(\mathbf{x})\cdot\mathbf{f}(\mathbf{x})$。
本章\emph{自治}指 $\mathbf{f}$ 不\textbf{显含时间} $t$（与“是否受控”无关）。
\end{note}

% ════════════════════════════════════════════════════════════════
\section{ODE、状态空间与机器人动力学}
\label{sec:ode-basic}

\paragraph{动机：机械臂的运动方程长什么样。}
设想最小的教学场景：一根绕固定轴转动的单连杆机械臂（末端带负载），在重力下受关节力矩 $\tau$ 驱动。用牛顿--欧拉或拉格朗日写出转动方程：
\begin{equation}
  I\ddot\theta+b\dot\theta+mg\ell\sin\theta=\tau,
  \label{eq:single-link}
\end{equation}
其中 $\theta$ 是关节角，$I$ 转动惯量，$b$ 黏性摩擦，$mg\ell\sin\theta$ 是重力恢复力矩。\textbf{这是一个二阶常微分方程}：未知量 $\theta(t)$ 是时间的函数，方程关联它的二阶导 $\ddot\theta$、一阶导 $\dot\theta$ 与它本身。给定初始角 $\theta(0)$、初始角速度 $\dot\theta(0)$ 与输入 $\tau(t)$，这根杆的轨迹是否被\emph{唯一}确定？拨动初始角 $0.01$ 弧度，一秒后末端差多少？设计反馈 $\tau=-k_p\theta-k_d\dot\theta$，杆会不会乖乖停在 $\theta=0$？这些问题——存在、唯一、敏感、稳定——\textbf{全部是 ODE 理论的问题}。这根单连杆只是最小例子；换成七自由度机械臂、四旋翼、双足，方程形式骤然复杂，但“它是一个 ODE”这件事不变。

\begin{insight}[ODE 是机器人在时间中如何变化的通用语法]
机器人学里几乎所有“连续时间发生的事”——关节怎么动、姿态怎么转、估计量怎么演化、最优轨迹满足什么方程——剥到最里层都是 $\dot{\mathbf{x}}=\mathbf{f}(t,\mathbf{x})$。ODE 不是机器人学的一个分支，而是描述机器人随时间变化的\emph{通用语法}。学控制、学估计、学仿真，本质上都是在学“如何分析、求解、利用这条 ODE 的不同侧面”。
\end{insight}

\paragraph{反面：抱着高阶形式不放，用不上任何工具。}
初学者常对着二阶方程 \cref{eq:single-link} 试图“直接积分两次”。对线性常系数情形偶尔行得通，但一旦出现 $\sin\theta$ 这种非线性项、或升到多自由度（$\mathbf{M}(\mathbf{q})$ 依赖 $\mathbf{q}$），就再没有“积两次”的套路。更糟的是，\textbf{几乎所有 ODE 的一般理论——存在唯一性、数值积分器、稳定性判据——都是对一阶系统 $\dot{\mathbf{x}}=\mathbf{f}(t,\mathbf{x})$ 陈述的}。抱着二阶形式不放，Picard--Lindelöf 定理、Lyapunov 函数全都用不上。反面教材：有人仿真单摆，写了个循环只更新 $\theta\mathrel{+}=\dot\theta\,h$ 却忘了同时更新 $\dot\theta$，因为他脑子里只有“$\theta$ 是状态”。养成“状态是 $(\theta,\dot\theta)$ 整体”的反射，这类错误根本不会发生。

\paragraph{历史：从 Newton 的“流数”到状态空间。}
ODE 与微积分同龄。Newton（1671 前后，《流数法》）研究的“流数方程”就是最早的微分方程；Leibniz 同期发展了 $dx/dt$ 记号。18 世纪 Euler、Lagrange、d'Alembert 把它推向力学，Lagrange 的分析力学（1788）系统地把质点系运动写成二阶微分方程组——这正是机器人欧拉--拉格朗日方程的祖先\cite{arnold1989mathematical}。把高阶方程统一化为一阶向量系统、并在\emph{相空间}里几何地看待解的轨迹，主要归功于 19 世纪末的 Poincaré（定性理论）与 20 世纪中叶控制论中的“状态空间方法”（Kalman 等，1960 前后）。今天说的“状态空间形式”$\dot{\mathbf{x}}=\mathbf{f}(t,\mathbf{x})$，是 Poincaré 几何直觉与控制论工程语言的合流。

\paragraph{理论：定义、分类与状态空间化。}

\begin{definition}[常微分方程与初值问题]\label{def:ode-ivp}
常微分方程是关于未知函数 $\mathbf{x}(t)$ 及其各阶导数的方程；出现的最高阶导数的阶称为方程的\emph{阶}。本章统一采用一阶向量形式的\textbf{初值问题}（Initial Value Problem, IVP）：
\begin{equation}
  \dot{\mathbf{x}}=\mathbf{f}(t,\mathbf{x}),\qquad \mathbf{x}(t_0)=\mathbf{x}_0,\qquad \mathbf{f}:U\subseteq\mathbb{R}\times\mathbb{R}^n\to\mathbb{R}^n,
  \label{eq:ivp}
\end{equation}
其中 $\dot{\mathbf{x}}=d\mathbf{x}/dt$，$\mathbf{x}_0$ 是初始条件，$U$ 是 $\mathbf{f}$ 的定义域（开集）。
\end{definition}

理解一个 ODE，先问它落在三组对立维度的哪一边——这决定了你能用哪些工具：

\begin{table}[H]\centering\small
\caption{ODE 的三组分类维度}
\begin{tabular}{llll}
\toprule
维度 & 一侧 & 另一侧 & 为何重要 \\
\midrule
阶数 & 一阶 $\dot{\mathbf{x}}=\mathbf{f}$ & 高阶 $\mathbf{x}^{(k)}=\mathbf{F}(\cdots)$ & 一般理论只对一阶陈述，高阶须先降阶 \\
线性性 & 线性 $\dot{\mathbf{x}}=\mathbf{A}(t)\mathbf{x}+\mathbf{g}$ & 非线性 $\dot{\mathbf{x}}=\mathbf{f}(t,\mathbf{x})$ & 线性可完全求解（\cref{sec:ode-linear}），非线性一般只能定性分析 \\
自治性 & 自治 $\dot{\mathbf{x}}=\mathbf{f}(\mathbf{x})$ & 非自治 $\dot{\mathbf{x}}=\mathbf{f}(t,\mathbf{x})$ & 自治系统的解定义相流，可画相图（\cref{sec:ode-phase}） \\
\bottomrule
\end{tabular}
\end{table}

\textbf{核心操作：高阶 $\to$ 一阶的状态扩张。} 任意 $k$ 阶标量方程 $x^{(k)}=F(t,x,\dot x,\dots,x^{(k-1)})$，通过引入新变量 $y_1=x,\,y_2=\dot x,\,\dots,\,y_k=x^{(k-1)}$，化为一阶系统
\begin{equation}
  \dot y_1=y_2,\quad \dot y_2=y_3,\quad \dots,\quad \dot y_{k-1}=y_k,\quad \dot y_k=F(t,y_1,\dots,y_k).
  \label{eq:order-reduction}
\end{equation}

\begin{insight}[降阶揭示了“状态”的真正含义]
降阶不是“技巧”，而是揭示了\emph{状态}的真正含义——要预测系统的未来，需要知道的全部当前信息，就是 $(x,\dot x,\dots,x^{(k-1)})$ 这一整组量。对机械系统，二阶方程意味着“位置 $+$ 速度”才构成完整状态；只知道位置预测不了未来，因为同一位置可以有任意速度，未来轨迹完全不同。这就是为什么\emph{相空间的维数是构型空间维数的两倍}。
\end{insight}

\textbf{自治系统与相流。} 当 $\mathbf{f}$ 不显含 $t$（自治系统 $\dot{\mathbf{x}}=\mathbf{f}(\mathbf{x})$），解 $\boldsymbol{\phi}(t,\mathbf{x}_0)$ 把“从 $\mathbf{x}_0$ 出发演化 $t$ 时间到哪”封装成映射族 $\boldsymbol{\phi}_t:\mathbb{R}^n\to\mathbb{R}^n$，称为\textbf{相流}（phase flow）。它满足群性质
\begin{equation}
  \boldsymbol{\phi}_{t+s}=\boldsymbol{\phi}_t\circ\boldsymbol{\phi}_s,\qquad \boldsymbol{\phi}_0=\mathrm{id},\qquad (\boldsymbol{\phi}_t)^{-1}=\boldsymbol{\phi}_{-t}.
  \label{eq:flow-group}
\end{equation}
直觉：先演化 $s$ 再演化 $t$，等于一口气演化 $t+s$——因为自治系统“规则不随时间改变”。可以记住：\emph{自治 ODE $=$ 相空间上的一个速度场，解 $=$ 顺着速度场漂流的轨迹}。这条群性质在微分流形语境（\cref{ch:lie} 的李群）将升级为“单参数微分同胚群”——\cref{eq:flow-group} 与李群指数映射的 $\exp((t+s)\mathbf{A})=\exp(t\mathbf{A})\exp(s\mathbf{A})$ 是同一件事的两个面孔。

\paragraph{机器人应用：欧拉--拉格朗日方程的状态空间形式。}
刚体机械臂的欧拉--拉格朗日方程（$\mathbf{q}\in\mathbb{R}^n$ 为关节坐标）：
\begin{equation}
  \mathbf{M}(\mathbf{q})\ddot{\mathbf{q}}+\mathbf{C}(\mathbf{q},\dot{\mathbf{q}})\dot{\mathbf{q}}+\mathbf{G}(\mathbf{q})=\boldsymbol{\tau},
  \label{eq:euler-lagrange}
\end{equation}
其中 $\mathbf{M}(\mathbf{q})\succ0$ 对称正定，$\mathbf{C}(\mathbf{q},\dot{\mathbf{q}})\dot{\mathbf{q}}$ 是科氏与离心力，$\mathbf{G}(\mathbf{q})$ 是重力，$\boldsymbol{\tau}$ 是关节力矩。这是 $n$ 维二阶系统。令状态 $\mathbf{x}=(\mathbf{q},\dot{\mathbf{q}})\in\mathbb{R}^{2n}$，因 $\mathbf{M}(\mathbf{q})$ 可逆，解出 $\ddot{\mathbf{q}}$ 得一阶状态空间形式：
\begin{equation}
  \dot{\mathbf{x}}=\begin{pmatrix}\dot{\mathbf{q}}\\[2pt]\mathbf{M}(\mathbf{q})^{-1}\bigl[\boldsymbol{\tau}-\mathbf{C}(\mathbf{q},\dot{\mathbf{q}})\dot{\mathbf{q}}-\mathbf{G}(\mathbf{q})\bigr]\end{pmatrix}=:\mathbf{f}(\mathbf{x},\boldsymbol{\tau}).
  \label{eq:robot-statespace}
\end{equation}
\textbf{这就是 MuJoCo、Drake、Isaac Sim、Pinocchio 等仿真器内部所有积分器的“被积对象”}。每一次 \texttt{sim.step()} 本质上都是用某个数值方法在小时间步 $h$ 内把这个 $\mathbf{f}$ 积一次。

\begin{insight}[理论-工程桥接：写成 $\mathbf{f}$ 与高效算出 $\mathbf{f}$ 是两件事]
注意 \cref{eq:robot-statespace} 里的 $\mathbf{M}(\mathbf{q})^{-1}$。仿真器不会真的去求 $\mathbf{M}$ 的逆（$O(n^3)$ 且数值差），而是用 Featherstone 的\emph{关节体算法}（Articulated Body Algorithm）在 $O(n)$ 时间内直接算出 $\ddot{\mathbf{q}}$。所以“写成 $\dot{\mathbf{x}}=\mathbf{f}$”是数学层面的概念统一，“高效算出 $\mathbf{f}$”是刚体动力学算法的工程任务——两者分工明确：ODE 理论管“解的性质”，动力学算法管“右端怎么快速求值”。
\end{insight}

\begin{pitfall}[概念误区：把“自治”当成“不受控”]
看到 $\dot{\mathbf{x}}=\mathbf{f}(\mathbf{x})$ 不含 $t$，就以为机器人不能受控、反馈系统都非自治——\emph{错}。“自治”指 $\mathbf{f}$ 不\textbf{显含时间 $t$}，与“有没有输入”无关。一个\emph{状态反馈}律 $\boldsymbol{\tau}=\mathbf{k}(\mathbf{x})$ 代回后得到的闭环系统 $\dot{\mathbf{x}}=\mathbf{f}(\mathbf{x},\mathbf{k}(\mathbf{x}))=:\tilde{\mathbf{f}}(\mathbf{x})$ 仍是\emph{自治}的！只有\emph{时变}输入 $\boldsymbol{\tau}(t)$（如开环时间轨迹）才让系统非自治。绝大多数稳定性分析（PD、Lyapunov、相图）针对的正是自治闭环系统——分清“开环时变轨迹”（非自治）与“状态反馈闭环”（自治），是用对这些工具的前提。
\end{pitfall}

\begin{pitfall}[思维陷阱：以为“写出了 $\mathbf{f}$ 就等于解决了问题”]
把动力学写成 $\dot{\mathbf{x}}=\mathbf{f}(\mathbf{x},\boldsymbol{\tau})$ 后就认为“建模完成”，直接交给数值积分而不问解是否存在、唯一、有界。$\mathbf{f}$ 只是“规则”，规则存在不代表“解存在且行为良好”——解的存在唯一性（\cref{sec:ode-picard}）、对扰动的敏感性、长时间行为是独立于“写出 $\mathbf{f}$”的数学问题。建模写出 $\mathbf{f}$ 后，例行检查 $\mathbf{f}$ 的正则性（是否 Lipschitz/$C^1$、增长是否有界），这是后续一切分析的前提。
\end{pitfall}

\begin{practice}
\begin{enumerate}
  \item 把单连杆方程 \cref{eq:single-link} 化为一阶状态空间形式 $\dot{\mathbf{x}}=\mathbf{f}(\mathbf{x},\tau)$，写出 $\mathbf{f}$ 的两个分量。令 $\tau\equiv0$、$b=0$，验证 $\mathbf{f}(\mathbf{x})$ 不显含 $t$（自治），并指出所有平衡点（$\mathbf{f}(\mathbf{x})=\mathbf{0}$ 的 $\mathbf{x}$）。它们对应单摆的什么物理位形？
  \item 含\emph{时滞}的系统 $\ddot\theta(t)=g(\theta(t),\theta(t-\tau_d))$（控制有 $\tau_d$ 秒延迟）的“完整状态”是什么？它还是有限维 ODE 吗？（提示：要预测未来需要知道哪一\emph{段}历史？这就是为什么延迟系统属于“泛函微分方程”而非普通 ODE——也是真机控制比仿真难的原因之一。）
  \item 把 PD$+$重力补偿律 $\boldsymbol{\tau}=-\mathbf{K}_p\mathbf{q}-\mathbf{K}_d\dot{\mathbf{q}}+\mathbf{G}(\mathbf{q})$ 代入 \cref{eq:robot-statespace}，写出闭环 $\dot{\mathbf{x}}=\tilde{\mathbf{f}}(\mathbf{x})$，验证它不显含 $t$（自治），并说明这为什么使我们能用相图与 Lyapunov 直接法（\cref{sec:ode-lyapunov}）分析它。
\end{enumerate}
\end{practice}

% ════════════════════════════════════════════════════════════════
\section{解的存在与唯一：Picard--Lindelöf 定理}
\label{sec:ode-picard}

\paragraph{动机：仿真器结果可复现吗。}
你在仿真器里设好初始状态 $\mathbf{x}_0$ 与力矩输入 $\boldsymbol{\tau}(t)$，点运行，得到一条轨迹。换台电脑、换个随机种子（但物理与输入不变），\emph{会不会得到同一条轨迹}？若今天得轨迹 A、明天得轨迹 B，那一切基于仿真的研究都无从谈起。这背后是最根本的数学问题：\textbf{给定 $\dot{\mathbf{x}}=\mathbf{f}(t,\mathbf{x})$、$\mathbf{x}(t_0)=\mathbf{x}_0$，解是否存在？是否唯一？} “存在”保证轨迹确实有；“唯一”保证它只有一条——这正是可复现性的数学定义。

\paragraph{反面：仅靠“$\mathbf{f}$ 连续”不够。}
一个自然的猜测是“$\mathbf{f}$ 连续就够了”。但连续\emph{只保证存在、不保证唯一}：经典反例 $\dot x=x^{2/3},\,x(0)=0$ 有无穷多条解（见 \cref{ex:nonunique}）。这意味着只要求连续，仿真器面对这种系统时“选哪条解”没有数学依据。我们需要比“连续”更强一点、又不太强的条件——恰到好处的就是 \textbf{Lipschitz 连续}。

\paragraph{历史：从 Cauchy 到 Picard 的逐次逼近。}
存在唯一性最早可追溯到 Cauchy（1820 年代，Cauchy--Lipschitz 定理）。Lipschitz（1876）引入了以他命名的条件简化论证。真正把证明变得优雅的是 Picard（1890）的\emph{逐次逼近法}：从一个初始猜测出发，反复代入积分方程得到越来越好的近似 $\mathbf{x}_0\to\mathbf{x}_1\to\mathbf{x}_2\to\cdots$，证明这串近似收敛到真解。Lindelöf（1894）把它形式化为现代版本。

\begin{insight}[Picard 逐次逼近 $=$ Banach 不动点定理在函数空间上的化身]
Picard 的逐次逼近 $\mathbf{x}_{n+1}(t)=\mathbf{x}_0+\int_{t_0}^t\mathbf{f}(s,\mathbf{x}_n(s))\,ds$ 不是孤立技巧——它正是 Banach 不动点定理在\emph{函数空间}上的化身。“求解 ODE”被翻译成“求一个积分算子 $T$ 的不动点”，而 Lipschitz 条件恰好让 $T$ 成为压缩映射。Picard 比 Banach 早约 30 年，相当于在具体问题上提前发现了抽象不动点定理的雏形。这是“抽象泛函分析如何统一具体分析问题”最漂亮的范例之一。
\end{insight}

\paragraph{理论：定理陈述与完整证明。}

\begin{theorem}[Picard--Lindelöf, 1890/1894]\label{thm:picard}
设 $\mathbf{f}$ 在闭矩形 $R=\{(t,\mathbf{x}):|t-t_0|\le a,\ \|\mathbf{x}-\mathbf{x}_0\|\le b\}$ 上\emph{连续}，且对 $\mathbf{x}$ 满足 \textbf{Lipschitz 条件}：存在常数 $L\ge0$ 使
\begin{equation}
  \|\mathbf{f}(t,\mathbf{x}_1)-\mathbf{f}(t,\mathbf{x}_2)\|\le L\,\|\mathbf{x}_1-\mathbf{x}_2\|,\qquad\forall (t,\mathbf{x}_1),(t,\mathbf{x}_2)\in R.
  \label{eq:lipschitz}
\end{equation}
记 $M=\sup_R\|\mathbf{f}\|$，$\delta=\min(a,\,b/M)$。则 IVP \cref{eq:ivp} 在 $[t_0-\delta,\,t_0+\delta]$ 上存在\emph{唯一}的 $C^1$ 解。
\end{theorem}

\begin{proof}
分四步：积分方程改写 $\to$ 算子自映 $\to$ 压缩 $\to$ 去掉技术限制。

\emph{第 (i) 步：化 IVP 为积分方程。} 若 $\mathbf{x}(t)$ 是 $C^1$ 解，对 $\dot{\mathbf{x}}=\mathbf{f}(t,\mathbf{x})$ 从 $t_0$ 到 $t$ 积分并用 $\mathbf{x}(t_0)=\mathbf{x}_0$：
\begin{equation}
  \mathbf{x}(t)=\mathbf{x}_0+\int_{t_0}^t\mathbf{f}\bigl(s,\mathbf{x}(s)\bigr)\,ds.
  \label{eq:integral-eq}
\end{equation}
反过来，若连续函数 $\mathbf{x}(t)$ 满足 \cref{eq:integral-eq}，右端可微（被积函数连续），求导得回 $\dot{\mathbf{x}}=\mathbf{f}(t,\mathbf{x})$ 且自动满足初值。所以“$C^1$ 解 IVP”$\iff$“连续解积分方程 \cref{eq:integral-eq}”。意义：把微分方程变成积分方程，而积分算子比微分算子温和（积分连续、求导不连续），便于做不动点。

\emph{第 (ii) 步：定义 Picard 算子并验证自映。} 取
\[
  X=\Bigl\{\mathbf{x}\in C\bigl([t_0-\delta,t_0+\delta],\mathbb{R}^n\bigr):\ \|\mathbf{x}(t)-\mathbf{x}_0\|\le b\ \forall t\Bigr\},
\]
配上上确界范数 $\|\mathbf{x}\|_\infty=\sup_t\|\mathbf{x}(t)\|$。$X$ 是完备空间 $C$ 中的闭球，故仍完备。定义 \textbf{Picard 算子} $(T\mathbf{x})(t):=\mathbf{x}_0+\int_{t_0}^t\mathbf{f}(s,\mathbf{x}(s))\,ds$。对 $\mathbf{x}\in X$、$|t-t_0|\le\delta$，
\[
  \|(T\mathbf{x})(t)-\mathbf{x}_0\|=\Bigl\|\int_{t_0}^t\mathbf{f}(s,\mathbf{x}(s))\,ds\Bigr\|\le M\,|t-t_0|\le M\delta\le M\cdot\frac{b}{M}=b,
\]
故 $T\mathbf{x}\in X$（这正是 $\delta\le b/M$ 这个取法的作用——保证迭代不跑出矩形 $R$，使 $\mathbf{f}$ 始终有定义）。

\emph{第 (iii) 步：验证 $T$ 是压缩映射。} 对 $\mathbf{x},\mathbf{y}\in X$，用 Lipschitz 条件：
\[
  \|(T\mathbf{x})(t)-(T\mathbf{y})(t)\|\le\int_{t_0}^t L\,\|\mathbf{x}(s)-\mathbf{y}(s)\|\,ds\le L|t-t_0|\,\|\mathbf{x}-\mathbf{y}\|_\infty.
\]
取上确界得 $\|T\mathbf{x}-T\mathbf{y}\|_\infty\le L\delta\,\|\mathbf{x}-\mathbf{y}\|_\infty$。若 $L\delta<1$，$T$ 是压缩，由 Banach 不动点定理有唯一不动点 $\mathbf{x}^*\in X$，即 IVP 的唯一解。

\emph{第 (iv) 步：去掉 $L\delta<1$ 的限制（加权范数技巧）。} 上面要求 $L\delta<1$，未必满足。解决办法不是缩小区间，而是\emph{换等价范数}让 $T$ 在整个区间上压缩。定义加权（Bielecki）范数 $\|\mathbf{x}\|_L:=\sup_{|t-t_0|\le\delta}e^{-2L|t-t_0|}\|\mathbf{x}(t)\|$，它与 $\|\cdot\|_\infty$ 等价。重做压缩估计（以 $t>t_0$ 为例）：
\[
  \|(T\mathbf{x})(t)-(T\mathbf{y})(t)\|\le L\int_{t_0}^t\|\mathbf{x}(s)-\mathbf{y}(s)\|\,ds\le L\int_{t_0}^t e^{2L(s-t_0)}ds\;\|\mathbf{x}-\mathbf{y}\|_L=\tfrac12\bigl(e^{2L(t-t_0)}-1\bigr)\|\mathbf{x}-\mathbf{y}\|_L.
\]
两边乘 $e^{-2L(t-t_0)}$ 取上确界：$\|T\mathbf{x}-T\mathbf{y}\|_L\le\tfrac12\|\mathbf{x}-\mathbf{y}\|_L$。\emph{压缩系数 $1/2<1$ 与 $\delta$ 无关}，于是无论 $\delta$ 多大，$T$ 在 $\|\cdot\|_L$ 下都是压缩，唯一不动点存在。
\end{proof}

\begin{insight}[加权范数：用指数因子驯服线性增长]
第 (iv) 步的加权范数是反复出现的深刻技巧——指数权 $e^{-2L|t-t_0|}$ 专门“压制”Lipschitz 常数 $L$ 带来的增长。它的精神与下文 Gronwall 估计、线性系统 $e^{t\mathbf{A}}$ 的作用完全一致（都是“用指数因子驯服线性增长”）。记住这个技巧：当一个迭代“差一点就压缩”时，往往可以靠重新加权救活。
\end{insight}

\begin{proposition}[$C^1$ 是 Lipschitz 的充分条件]\label{prop:c1-lipschitz}
若 $\mathbf{f}$ 关于 $\mathbf{x}$ 连续可微且 $\partial\mathbf{f}/\partial\mathbf{x}$ 在 $R$ 上有界，则 $\mathbf{f}$ 对 $\mathbf{x}$ Lipschitz，$L=\sup_R\|\partial\mathbf{f}/\partial\mathbf{x}\|$（由中值定理）。所以实践中只要 $\mathbf{f}\in C^1$，局部就自动 Lipschitz。
\end{proposition}

\textbf{局部 vs 全局。} \cref{thm:picard} 只保证 $[t_0-\delta,t_0+\delta]$ 上的\emph{局部}存在唯一；$\delta$ 由粗糙的全局界 $M=\sup\|\mathbf{f}\|$ 决定，常远小于解的真实寿命。解能否延伸到更长时间，是\emph{延拓理论}的主题：要么全局存在，要么\emph{有限时间爆破}（解在某 $t^*<\infty$ 处趋于无穷，例如 $\dot x=x^2,\,x(0)=1$ 的解 $x(t)=1/(1-t)$ 在 $t\to1^-$ 爆破——这正是“积分器发散”最常见的数学本质）。但唯一性是“局部到哪、唯一到哪”：两条解只要在某点重合，必处处重合。

\paragraph{机器人应用：仿真可复现性的数学前提。}
机器人动力学的状态空间右端 \cref{eq:robot-statespace} 在任意有界状态集上是 $C^1$ 的——因为 $\mathbf{M}(\mathbf{q})\succ0$ 连续可逆，$\mathbf{M}^{-1}$ 是 $\mathbf{M}$ 元素的光滑函数（Cramer 法则 $+$ 行列式非零），$\mathbf{C},\mathbf{G}$ 由动力学的光滑函数构成。$C^1\Rightarrow$ 局部 Lipschitz（\cref{prop:c1-lipschitz}）$\Rightarrow$ 给定 $\mathbf{x}_0$ 与连续 $\boldsymbol{\tau}(t)$，\textbf{机械臂轨迹唯一存在}。

\begin{insight}[理论-工程桥接：可复现性的数学保证及其崩塌]
只要 $\boldsymbol{\tau}(t)$ 连续、动力学光滑，同样的 $\mathbf{x}_0$ 与 $\boldsymbol{\tau}(t)$ 必给出同一条轨迹——与电脑、随机种子无关。反之，当控制律含 $\operatorname{sign}(\cdot)$（滑模）这种\emph{不连续}项时，Lipschitz 失效，唯一性可能崩塌：仿真器选哪条解取决于离散方式，不同步长 / 求解器给出不同结果，“可复现”不再自动成立。从建模层面恢复唯一性（Filippov 解、松弛化软接触使 $\mathbf{f}$ 重新 Lipschitz）比 debug 代码更对症。
\end{insight}

\begin{example}[连续不保证唯一：$\dot x=x^{2/3}$]\label{ex:nonunique}
标量 IVP $\dot x=x^{2/3},\,x(0)=0$，右端 $f(x)=x^{2/3}$ \emph{连续}但在 $x=0$ 处\emph{不 Lipschitz}（导数 $\tfrac23 x^{-1/3}$ 在 $0$ 爆掉）。它至少有两条解：恒零解 $x(t)\equiv0$，以及 $x(t)=(t/3)^3=t^3/27$（验证：$\dot x=\tfrac{t^2}{9}$，而 $x^{2/3}=(t^3/27)^{2/3}=t^2/9$，两端相等）。事实上，对任意 $c\ge0$，函数“在 $[0,c]$ 上取 $0$、之后取 $((t-c)/3)^3$”都是解——\emph{无穷多条}。这正是 Lipschitz 失效导致唯一性崩塌的最小例子，也是非光滑机器人动力学“宽进严出”的数学根源。
\end{example}

\begin{pitfall}[概念误区：把 Lipschitz 当成“可微”或“连续”]
三个条件按\emph{函数类}的真包含排成 $C^1\subsetneq$ Lipschitz $\subsetneq$ 连续（局部意义）：每个（导数有界的）$C^1$ 函数都 Lipschitz、每个 Lipschitz 函数都连续，反向均不成立，故越靠左条件越强、函数类越小。唯一性要的恰是中间的 Lipschitz——\emph{比连续强}（排除 $x^{2/3}$），\emph{比 $C^1$ 弱}（容纳含 $|x|$ 的项，它处处 Lipschitz 但在 $0$ 不可微）。判断唯一性查 Lipschitz；$\mathbf{f}\in C^1$ 且导数有界是\emph{充分}条件但不必要。含绝对值 / 分段线性的项（如饱和函数 $\operatorname{sat}$）通常 Lipschitz，唯一性仍成立；只有 $\operatorname{sign}$ 这种\emph{跳变}才真正破坏 Lipschitz。
\end{pitfall}

\begin{pitfall}[概念误区：以为 $\delta$ 就是解的真实寿命]
算出 $\delta=\min(a,b/M)$ 后认为“解只活到 $t_0+\delta$”——错。$\delta$ 是\emph{保证局部存在的下界}（用了粗糙的全局界 $M$），靠延拓可把解一步步接到更长区间，直到真正爆破或逸出。把 $\delta$ 理解为“至少能活这么久”；真实最大存在区间由延拓定理刻画：要么全局存在，要么有限时间爆破。
\end{pitfall}

\begin{practice}
\begin{enumerate}
  \item 完整复述 \cref{thm:picard} 的四步证明，不看讲义。特别地，对 $t<t_0$ 的情形（区间左半边）重做第 (iv) 步加权范数的压缩估计，验证压缩系数仍为 $1/2$。
  \item 对标量 IVP $\dot x=2x,\,x(0)=1$，从 $x_0(t)\equiv1$ 出发做三步 Picard 迭代算 $x_1,x_2,x_3$。观察它们是否在向 $e^{2t}$ 的 Taylor 级数逼近，写出第 $n$ 步通项并说明 $\lim_n x_n=e^{2t}$。
  \item 考虑饱和反馈 $\tau=-\operatorname{sat}(k_p\theta+k_d\dot\theta)$，$\operatorname{sat}(s)=\max(-1,\min(1,s))$ 连续且分段线性。代入单摆后的闭环 $\dot{\mathbf{x}}=\tilde{\mathbf{f}}(\mathbf{x})$ 是否 Lipschitz？解是否唯一？把 $\operatorname{sat}$ 换成 $\operatorname{sign}$ 呢？分别论证，并说明这对“为什么饱和控制比纯继电（bang-bang）控制在仿真中更乖”的工程含义。
\end{enumerate}
\end{practice}

% ════════════════════════════════════════════════════════════════
\section{线性系统：基本解矩阵与常数变易}
\label{sec:ode-linear}

\paragraph{动机：唯一能“完全解出来”的一大类。}
非线性 ODE 一般没有闭式解——只能定性分析、数值积分或证稳定性。但有一大类例外：\textbf{线性系统 $\dot{\mathbf{x}}=\mathbf{A}(t)\mathbf{x}+\mathbf{g}(t)$ 可以完全求解}，解有显式公式。这不只是便利——它是整个线性控制理论（LQR、Kalman 滤波、极点配置）的根基，而这些理论又通过“在平衡点附近线性化”成为非线性机器人控制的局部工具。所以\emph{理解线性系统 $=$ 理解扰动如何传播 $=$ 理解局部稳定性}。

\paragraph{反面：不会线性系统会怎样。}
跳过线性理论，你将无法理解控制工程的“通用语言”：所有 LTI（线性时不变）系统的能控性、能观性、稳定性、最优控制都建立在 $\dot{\mathbf{x}}=\mathbf{A}\mathbf{x}+\mathbf{B}\mathbf{u}$ 之上。看到 LQR 的 Riccati 方程、Kalman 滤波的协方差传播、状态转移矩阵 $\boldsymbol{\Phi}(t,t_0)$，你会一头雾水——它们全是本节“基本解矩阵 $+$ 常数变易”的直接产物\cite{chen1999linear}。

\paragraph{历史。}
线性 ODE 的系统理论可追溯到 18 世纪（Euler、Lagrange 的常数变易法），矩阵形式的处理随线性代数成熟于 19--20 世纪。“状态转移矩阵”“基本解矩阵”等术语在 20 世纪中叶随控制论定型。常数变易公式（variation of parameters）是 Lagrange 的发明——一个“把齐次解的常数换成时间函数”的精妙技巧。

\paragraph{理论：齐次系统、基本解矩阵与常数变易。}

\textbf{齐次系统 $\dot{\mathbf{x}}=\mathbf{A}(t)\mathbf{x}$。} 解集是 $n$ 维线性空间（叠加原理：解的线性组合仍是解）。取 $n$ 个线性无关解作列，组成\textbf{基本解矩阵}（fundamental matrix）$\boldsymbol{\Phi}(t)\in\mathbb{R}^{n\times n}$，满足
\begin{equation}
  \dot{\boldsymbol{\Phi}}=\mathbf{A}(t)\boldsymbol{\Phi},\qquad \det\boldsymbol{\Phi}(t)\neq0\ \ \forall t.
  \label{eq:fund-matrix}
\end{equation}
\textbf{状态转移矩阵} $\boldsymbol{\Phi}(t,t_0):=\boldsymbol{\Phi}(t)\boldsymbol{\Phi}(t_0)^{-1}$ 满足 $\boldsymbol{\Phi}(t_0,t_0)=\mathbf{I}$，把初值映到 $t$ 时刻：$\mathbf{x}(t)=\boldsymbol{\Phi}(t,t_0)\mathbf{x}_0$。

\begin{proposition}[Liouville 公式：相空间体积的演化]\label{prop:liouville}
$W(t):=\det\boldsymbol{\Phi}(t)$ 满足 $\dot W=\operatorname{tr}(\mathbf{A}(t))\,W$，故
\begin{equation}
  W(t)=W(t_0)\exp\Bigl(\int_{t_0}^t\operatorname{tr}\mathbf{A}(s)\,ds\Bigr).
  \label{eq:liouville}
\end{equation}
意义：相空间体积的变化率由 $\operatorname{tr}\mathbf{A}$ 决定——$\operatorname{tr}\mathbf{A}<0$ 体积收缩（耗散），$\operatorname{tr}\mathbf{A}=0$ 体积守恒（保守 / 哈密顿系统），$\operatorname{tr}\mathbf{A}>0$ 体积膨胀。
\end{proposition}
\begin{proof}
用 Jacobi 公式 $\frac{d}{dt}\det\boldsymbol{\Phi}=\det\boldsymbol{\Phi}\cdot\operatorname{tr}(\boldsymbol{\Phi}^{-1}\dot{\boldsymbol{\Phi}})$，代入 $\dot{\boldsymbol{\Phi}}=\mathbf{A}\boldsymbol{\Phi}$ 得 $\operatorname{tr}(\boldsymbol{\Phi}^{-1}\mathbf{A}\boldsymbol{\Phi})=\operatorname{tr}\mathbf{A}$（迹在相似变换下不变），即 $\dot W=(\operatorname{tr}\mathbf{A})W$，积分得 \cref{eq:liouville}。
\end{proof}

\textbf{非齐次系统：常数变易公式。} 对 $\dot{\mathbf{x}}=\mathbf{A}(t)\mathbf{x}+\mathbf{g}(t)$，设解形如 $\mathbf{x}(t)=\boldsymbol{\Phi}(t)\mathbf{c}(t)$（把齐次解 $\boldsymbol{\Phi}\mathbf{c}$ 的常向量 $\mathbf{c}$ 换成 $\mathbf{c}(t)$），代入用 $\dot{\boldsymbol{\Phi}}=\mathbf{A}\boldsymbol{\Phi}$ 推出 $\boldsymbol{\Phi}\dot{\mathbf{c}}=\mathbf{g}$，积分得
\begin{equation}
  \boxed{\,\mathbf{x}(t)=\boldsymbol{\Phi}(t,t_0)\,\mathbf{x}_0+\int_{t_0}^t\boldsymbol{\Phi}(t,s)\,\mathbf{g}(s)\,ds.\,}
  \label{eq:variation-constants}
\end{equation}
\textbf{结构解读}：第一项是初值的\emph{自由演化}（零输入响应），第二项是输入 $\mathbf{g}(s)$ 在各时刻被“传播”到 $t$ 的叠加（零状态响应）——每个时刻 $s$ 注入的 $\mathbf{g}(s)$ 被状态转移矩阵 $\boldsymbol{\Phi}(t,s)$ 送到 $t$，再积分累加。当 $\mathbf{A}$ 为常数，$\boldsymbol{\Phi}(t,t_0)=e^{(t-t_0)\mathbf{A}}$（矩阵指数，\cref{sec:ode-expm}），公式变为 $\mathbf{x}(t)=e^{(t-t_0)\mathbf{A}}\mathbf{x}_0+\int_{t_0}^t e^{(t-s)\mathbf{A}}\mathbf{g}(s)\,ds$。

\begin{insight}[常数变易公式是“线性 $=$ 叠加”的终极表达]
任何线性系统的响应都能分解为“零输入响应”（初值自由演化）$+$“零状态响应”（输入的卷积积分）。这个分解在频域就是传递函数、在时域就是脉冲响应卷积、在控制里就是“自由 $+$ 强迫”。$\boldsymbol{\Phi}(t,s)$ 扮演格林函数 / 脉冲响应的角色——它是线性系统全部信息的载体。
\end{insight}

\paragraph{机器人应用：所有 LTI 控制理论的起点。}
机器人系统在平衡点 $\mathbf{x}^*$（或标称轨迹）附近线性化：$\delta\dot{\mathbf{x}}=\mathbf{A}\,\delta\mathbf{x}+\mathbf{B}\,\delta\mathbf{u}$，$\mathbf{A}=D_\mathbf{x}\mathbf{f}|_{\mathbf{x}^*}$、$\mathbf{B}=D_\mathbf{u}\mathbf{f}|_{\mathbf{x}^*}$。这是所有线性控制理论的入口：能控矩阵 $\mathcal{C}=[\mathbf{B},\mathbf{A}\mathbf{B},\dots,\mathbf{A}^{n-1}\mathbf{B}]$ 满秩 $\iff$ 可任意配置极点；设计 $\mathbf{u}=-\mathbf{K}\mathbf{x}$ 使闭环 $\mathbf{A}-\mathbf{B}\mathbf{K}$ 的特征值在期望位置（LQR 通过 Riccati 方程，\cref{sec:ode-lyapeq}）；观测器误差动力学 $\dot{\mathbf{e}}=(\mathbf{A}-\mathbf{L}\mathbf{C})\mathbf{e}$ 是线性系统，协方差传播用状态转移矩阵——这正是连续时间 Kalman 滤波（\cref{ch:eskf}）的连续模型骨架。

\begin{pitfall}[概念误区：时变系统误用 $e^{\int\mathbf{A}\,ds}$]
对时变 $\mathbf{A}(t)$ 写解为 $\exp(\int_{t_0}^t\mathbf{A}(s)\,ds)$——除非 $\mathbf{A}(t)$ 与其积分\emph{可交换}，否则\textbf{错误}。原因：$\frac{d}{dt}\exp(\int\mathbf{A})=\mathbf{A}\exp(\int\mathbf{A})$ 仅当 $\mathbf{A}(t)$ 与 $\int\mathbf{A}$ 交换才成立（一般不成立，因矩阵不交换）。时变系统一般无闭式 $\boldsymbol{\Phi}$，须数值积分 $\dot{\boldsymbol{\Phi}}=\mathbf{A}(t)\boldsymbol{\Phi}$；只有 $\mathbf{A}$ 常数或处处交换时才有 $\boldsymbol{\Phi}=e^{\int\mathbf{A}}$。另：基本解矩阵\emph{不唯一}（右乘任意可逆常矩阵仍是），唯一的是状态转移矩阵 $\boldsymbol{\Phi}(t,t_0)$。
\end{pitfall}

\begin{practice}
\begin{enumerate}
  \item 用常数变易法完整推导 \cref{eq:variation-constants}：设 $\mathbf{x}=\boldsymbol{\Phi}(t)\mathbf{c}(t)$ 代入 $\dot{\mathbf{x}}=\mathbf{A}\mathbf{x}+\mathbf{g}$，利用 $\dot{\boldsymbol{\Phi}}=\mathbf{A}\boldsymbol{\Phi}$ 推出 $\boldsymbol{\Phi}\dot{\mathbf{c}}=\mathbf{g}$、$\dot{\mathbf{c}}=\boldsymbol{\Phi}^{-1}\mathbf{g}$，积分得公式。每步说清用了什么（叠加、$\boldsymbol{\Phi}$ 可逆）。
  \item 证明 Liouville 公式 \cref{eq:liouville}，并解释为什么 $\operatorname{tr}\mathbf{A}=0$ 对应相空间体积守恒，联系保守（哈密顿）系统。
  \item 机器人在标称轨迹（非平衡）附近线性化得 $\delta\dot{\mathbf{x}}=\mathbf{A}(t)\delta\mathbf{x}$，$\mathbf{A}(t)$ 随时间变化。为什么这种情形不能简单用单个 $e^{t\mathbf{A}}$ 分析稳定性？（提示：联系上面那个陷阱；这也是周期轨道 Floquet 理论要专门处理的情形。）
\end{enumerate}
\end{practice}

% ════════════════════════════════════════════════════════════════
\section{矩阵指数与 \texorpdfstring{$\SO(3)$}{SO(3)} 指数映射}
\label{sec:ode-expm}

\paragraph{动机：线性系统的万能解算子与姿态积分。}
常系数线性系统 $\dot{\mathbf{x}}=\mathbf{A}\mathbf{x}$ 的解是 $\mathbf{x}(t)=e^{t\mathbf{A}}\mathbf{x}_0$——矩阵指数 $e^{t\mathbf{A}}$ 是线性系统的“万能解算子”。而它的另一副面孔更贴近机器人：四旋翼、IMU、SLAM 中任何三维旋转都要回答\emph{给定角速度 $\boldsymbol{\omega}(t)$，姿态 $\mathbf{R}(t)\in\SO(3)$ 怎么积分}。姿态运动学 $\dot{\mathbf{R}}=\mathbf{R}[\boldsymbol{\omega}]_\times$ 是一个在\emph{矩阵群}上的 ODE，最干净的解法正是矩阵指数——它把“在李代数（角速度）里走一步”映成“李群（旋转）里的真旋转”。本节把 $e^{t\mathbf{A}}$ 讲透，是通向 \cref{ch:lie} 李群与下游 SLAM 的关键一跃。

\paragraph{反面：用幂级数硬算或用欧拉角会怎样。}
朴素地想，$e^{t\mathbf{A}}=\sum t^k\mathbf{A}^k/k!$ 是无穷级数，似乎得截断求和——但截断有误差，对大 $\|t\mathbf{A}\|$ 收敛慢、数值差（著名的“计算矩阵指数的十九种可疑方法”，Moler--Van Loan）。而对姿态，用欧拉角积分会在特定姿态遇\emph{万向锁}（gimbal lock）：两个旋转轴重合、丢一个自由度，积分公式含 $1/\cos(\text{pitch})$ 在 $\pm90^\circ$ 爆掉。矩阵指数的妙处在于：对结构特殊的矩阵（可对角化、Jordan 形、反对称），无穷级数能\emph{闭式求和}成有限表达式——对 $\SO(3)$ 就是著名的罗德里格斯公式。

\paragraph{理论：定义、性质与计算。}

\begin{definition}[矩阵指数]\label{def:matrix-exp}
对方阵 $\mathbf{A}\in\mathbb{R}^{n\times n}$，$e^{t\mathbf{A}}:=\sum_{k=0}^\infty\frac{t^k\mathbf{A}^k}{k!}$。该级数在任意矩阵范数下\emph{绝对收敛}（因 $\sum\|t\mathbf{A}\|^k/k!=e^{\|t\mathbf{A}\|}<\infty$），故良定义。
\end{definition}

\begin{proposition}[矩阵指数的基本性质]\label{prop:exp-props}
\begin{enumerate}
  \item \emph{导数}：$\frac{d}{dt}e^{t\mathbf{A}}=\mathbf{A}e^{t\mathbf{A}}=e^{t\mathbf{A}}\mathbf{A}$，故 $e^{t\mathbf{A}}$ 是 $\dot{\boldsymbol{\Phi}}=\mathbf{A}\boldsymbol{\Phi},\,\boldsymbol{\Phi}(0)=\mathbf{I}$ 的解（这就是 $\dot{\mathbf{x}}=\mathbf{A}\mathbf{x}$ 解为 $e^{t\mathbf{A}}\mathbf{x}_0$ 的根据）；
  \item \emph{半群}：$e^{(s+t)\mathbf{A}}=e^{s\mathbf{A}}e^{t\mathbf{A}}$（对同一 $\mathbf{A}$ 总成立，对应相流群性质 \cref{eq:flow-group}）；
  \item \emph{和}：$\mathbf{A}\mathbf{B}=\mathbf{B}\mathbf{A}\Rightarrow e^{\mathbf{A}+\mathbf{B}}=e^{\mathbf{A}}e^{\mathbf{B}}$——\textbf{仅当 $\mathbf{A},\mathbf{B}$ 交换！}一般不成立；
  \item \emph{逆}：$(e^{\mathbf{A}})^{-1}=e^{-\mathbf{A}}$，故 $e^{\mathbf{A}}$ 总可逆；
  \item \emph{行列式}：$\det e^{\mathbf{A}}=e^{\operatorname{tr}\mathbf{A}}$（连接 Liouville 公式 \cref{eq:liouville}）；
  \item \emph{谱映射}：$\sigma(e^{t\mathbf{A}})=\{e^{t\lambda}:\lambda\in\sigma(\mathbf{A})\}$——特征值跟着指数化。
\end{enumerate}
\end{proposition}

\begin{pitfall}[概念误区：误以为 $e^{\mathbf{A}+\mathbf{B}}=e^{\mathbf{A}}e^{\mathbf{B}}$ 总成立]
性质 (iii) 是最大陷阱。$e^{\mathbf{A}+\mathbf{B}}\neq e^{\mathbf{A}}e^{\mathbf{B}}$ 当 $\mathbf{A},\mathbf{B}$ 不交换——这正是\emph{旋转不可交换}（先转 X 再转 Y $\neq$ 先转 Y 再转 X）的代数根源，也是 BCH 公式（\cref{ch:lie}）要修正的地方。合成旋转 / 平移用 $e^{\mathbf{A}}e^{\mathbf{B}}$（按顺序，群乘法），不要先加指数。
\end{pitfall}

\textbf{计算方法。}
\emph{(1) 对角化}：$\mathbf{A}=\mathbf{P}\mathbf{D}\mathbf{P}^{-1}$（$\mathbf{D}=\operatorname{diag}(\lambda_i)$）$\Rightarrow e^{t\mathbf{A}}=\mathbf{P}\operatorname{diag}(e^{t\lambda_i})\mathbf{P}^{-1}$。最简单，但仅适用于可对角化的 $\mathbf{A}$。
\emph{(2) Jordan 形}（连接 \cref{ch:linalg}）：一般 $\mathbf{A}=\mathbf{P}\mathbf{J}\mathbf{P}^{-1}$，$\mathbf{J}=\mathbf{D}+\mathbf{N}$（$\mathbf{D}$ 对角、$\mathbf{N}$ 幂零、$\mathbf{D}\mathbf{N}=\mathbf{N}\mathbf{D}$）。由性质 (iii)（$\mathbf{D},\mathbf{N}$ 交换），
\begin{equation}
  e^{t\mathbf{J}}=e^{t\mathbf{D}}e^{t\mathbf{N}},\qquad e^{t\mathbf{N}}=\sum_{k=0}^{n-1}\frac{t^k\mathbf{N}^k}{k!}\quad(\text{有限和，因 }\mathbf{N}^n=\mathbf{0}).
  \label{eq:expm-jordan}
\end{equation}
\emph{幂零部分让无穷级数退化为多项式}——这就是 Jordan 形的威力：把 $e^{t\mathbf{A}}$ 化成“指数 $\times$ 多项式”的闭式。重特征值产生 $t^k e^{t\lambda}$ 形式的项（对应临界阻尼、重根响应）。
\emph{(3) 其他}：Cayley--Hamilton（$e^{t\mathbf{A}}$ 是 $\mathbf{A}$ 的 $\le n-1$ 次多项式）、Laplace 变换 $\mathcal{L}\{e^{t\mathbf{A}}\}=(s\mathbf{I}-\mathbf{A})^{-1}$、数值上的 scaling-and-squaring $+$ Padé（\texttt{scipy.\allowbreak linalg.\allowbreak expm}、Eigen \texttt{.exp()}）。

\paragraph{招牌闭式：罗德里格斯公式（$\so(3)$ 上的指数映射）。}
设 $\boldsymbol{\omega}\in\mathbb{R}^3$（旋转轴 $\times$ 角速度），其反对称矩阵 $[\boldsymbol{\omega}]_\times\in\so(3)$ 满足 $[\boldsymbol{\omega}]_\times\mathbf{v}=\boldsymbol{\omega}\times\mathbf{v}$。记 $\theta=\|\boldsymbol{\omega}\|$、单位轴 $\hat{\boldsymbol{\omega}}=\boldsymbol{\omega}/\theta$。

\begin{theorem}[罗德里格斯公式]\label{thm:rodrigues}
\begin{equation}
  e^{[\boldsymbol{\omega}]_\times}=\mathbf{I}+\sin\theta\,[\hat{\boldsymbol{\omega}}]_\times+(1-\cos\theta)\,[\hat{\boldsymbol{\omega}}]_\times^2\ \in\SO(3).
  \label{eq:rodrigues-ode}
\end{equation}
\end{theorem}
\begin{proof}
关键恒等式 $[\hat{\boldsymbol{\omega}}]_\times^3=-[\hat{\boldsymbol{\omega}}]_\times$（由 $\hat{\boldsymbol{\omega}}\times(\hat{\boldsymbol{\omega}}\times(\hat{\boldsymbol{\omega}}\times\mathbf{v}))=-\hat{\boldsymbol{\omega}}\times\mathbf{v}$，用 $\|\hat{\boldsymbol{\omega}}\|=1$）。于是 $[\hat{\boldsymbol{\omega}}]_\times$ 的幂循环：$[\hat{\boldsymbol{\omega}}]_\times^3=-[\hat{\boldsymbol{\omega}}]_\times$，$[\hat{\boldsymbol{\omega}}]_\times^4=-[\hat{\boldsymbol{\omega}}]_\times^2$，$\dots$ 代入级数 $e^{\theta[\hat{\boldsymbol{\omega}}]_\times}=\sum_k\frac{\theta^k}{k!}[\hat{\boldsymbol{\omega}}]_\times^k$，按奇次幂（含 $[\hat{\boldsymbol{\omega}}]_\times$）与偶次幂（含 $[\hat{\boldsymbol{\omega}}]_\times^2$）分组：
\[
  \text{奇次和}=\Bigl(\theta-\tfrac{\theta^3}{3!}+\tfrac{\theta^5}{5!}-\cdots\Bigr)[\hat{\boldsymbol{\omega}}]_\times=\sin\theta\,[\hat{\boldsymbol{\omega}}]_\times,\quad
  \text{偶次和}(k\ge2)=\Bigl(\tfrac{\theta^2}{2!}-\tfrac{\theta^4}{4!}+\cdots\Bigr)[\hat{\boldsymbol{\omega}}]_\times^2=(1-\cos\theta)[\hat{\boldsymbol{\omega}}]_\times^2.
\]
加上 $k=0$ 项 $\mathbf{I}$，得 \cref{eq:rodrigues-ode}。三项有限和，无需无穷级数、无万向锁。
\end{proof}
\noindent 这个公式与 \cref{ch:rigid_body} 由轴角求旋转矩阵的罗德里格斯公式\emph{一模一样}——那里从几何（向量绕轴旋转）得到它，这里从“$e^{t\mathbf{A}}$ 解矩阵 ODE $\dot{\mathbf{R}}=\mathbf{R}[\boldsymbol{\omega}]_\times$”的代数角度重新得到，两条路殊途同归。\cref{ch:lie} 会把它确立为 $\so(3)\to\SO(3)$ 的指数映射。

\begin{insight}[Euler 公式与 Rodrigues 公式是同一现象的两个面孔]
罗德里格斯公式是“$e^{t\mathbf{A}}$ 在 $\mathbf{A}^3=-\mathbf{A}$ 这种循环结构下闭式求和”的特例——与“$e^{i\theta}=\cos\theta+i\sin\theta$ 来自 $i^2=-1$ 的循环”一模一样。$[\hat{\boldsymbol{\omega}}]_\times$ 扮演了“矩阵版虚数单位”的角色（满足 $[\hat{\boldsymbol{\omega}}]_\times^3=-[\hat{\boldsymbol{\omega}}]_\times$，类比 $i^3=-i$）。这就是为什么二维旋转用复数、三维旋转用旋转矩阵 / 四元数：它们都是“指数映射把代数转成几何”。
\end{insight}

\paragraph{机器人应用：姿态积分与流形积分器。}
姿态运动学 $\dot{\mathbf{R}}=\mathbf{R}[\boldsymbol{\omega}]_\times$ 在小时间步 $\Delta t$ 内（$\boldsymbol{\omega}$ 近似常值）的解是
\begin{equation}
  \mathbf{R}(t+\Delta t)=\mathbf{R}(t)\,\exp\bigl([\boldsymbol{\omega}]_\times\Delta t\bigr),
  \label{eq:attitude-integration}
\end{equation}
用 \cref{eq:rodrigues-ode} 闭式展开（$\theta'=\|\boldsymbol{\omega}\|\Delta t$）。\textbf{结果严格落在 $\SO(3)$ 上}（$\exp$ 把 $\so(3)$ 映进 $\SO(3)$），无需事后重正交化——这是流形积分相对于“欧拉角 / 朴素积分 $+$ 归一化”的根本优势。把两个关键帧之间的高频 IMU 测量\emph{预积分}成一个相对运动约束，核心运算正是 $\SO(3)$ 上的指数 / 对数映射与右雅可比——这是 VINS、ORB-SLAM3、GTSAM 等 VIO 系统的数学内核，详见 \cref{ch:lie} 与 \cref{ch:eskf}。

\begin{pitfall}[思维陷阱：用截断幂级数数值算 $e^{t\mathbf{A}}$，或用欧拉角积分姿态]
两个常见错误：(a) 直接 $\sum_{k=0}^{K}t^k\mathbf{A}^k/k!$ 截断求和——$\|t\mathbf{A}\|$ 大时收敛慢、相减相消、舍入误差大（Moler--Van Loan 警示）；正确做法是 scaling-and-squaring $+$ Padé，对反对称矩阵直接用 Rodrigues 闭式（最快最稳）。(b) 用 roll-pitch-yaw 欧拉角积分 $\dot{\mathbf{R}}$，在 pitch $=\pm90^\circ$ 处奇异发散；正确做法是 $\mathbf{R}\leftarrow\mathbf{R}\exp([\boldsymbol{\omega}]_\times\Delta t)$ 或四元数积分，全程无奇异。
\end{pitfall}

\begin{practice}
\begin{enumerate}
  \item 验证 $[\hat{\boldsymbol{\omega}}]_\times^3=-[\hat{\boldsymbol{\omega}}]_\times$（$\|\hat{\boldsymbol{\omega}}\|=1$），独立推导罗德里格斯公式（奇偶幂分组求和）。再推导 $\SO(3)$ 对数映射 $\theta=\arccos\frac{\operatorname{tr}\mathbf{R}-1}{2}$ 并指出它在 $\theta\to0,\pi$ 处的数值困难。
  \item 对二维旋转生成元 $\mathbf{J}=\bigl(\begin{smallmatrix}0&-1\\1&0\end{smallmatrix}\bigr)$（满足 $\mathbf{J}^2=-\mathbf{I}$），用幂级数证明 $e^{\theta\mathbf{J}}=\bigl(\begin{smallmatrix}\cos\theta&-\sin\theta\\\sin\theta&\cos\theta\end{smallmatrix}\bigr)$。对比 Euler 公式说明 $\mathbf{J}$ 与 $i$ 的对应——这是上面“矩阵版虚数单位”洞察的二维版。
  \item 论证：朴素积分 $\mathbf{R}\leftarrow\mathbf{R}+\mathbf{R}[\boldsymbol{\omega}]_\times\Delta t$（一阶 Euler）多步后 $\mathbf{R}$ 偏离正交（$\det$ 偏离 $1$、列不再正交），需周期重正交化；而 \cref{eq:attitude-integration} 始终在 $\SO(3)$ 上。为什么说指数映射积分是“内蕴”的、朴素积分是“外蕴 $+$ 投影”的？
\end{enumerate}
\end{practice}

% ════════════════════════════════════════════════════════════════
\section{相平面、不动点与线性化稳定性}
\label{sec:ode-phase}

\paragraph{动机：非线性系统局部“长什么样”。}
倒立摆、PD 控制的机械臂、双足站立平衡——这些都是非线性系统，一般解不出闭式。但我们最关心的往往是\emph{平衡点附近的局部行为}：摆在竖直位置稳不稳？PD 下机械臂怎么趋近目标？这些“局部定性行为”能不能不解方程就看出来？答案是能——通过\textbf{在平衡点处线性化}，看线性化矩阵 $\mathbf{A}=D\mathbf{f}(\mathbf{x}^*)$ 的特征值。但有个深刻问题：线性化是近似，\emph{凭什么相信它对原非线性系统的定性行为可靠}？Hartman--Grobman 定理给出精确回答。

\paragraph{反面：盲目相信或盲目不信线性化。}
两个极端都有害。\emph{盲目相信}：以为线性化在任何平衡点都可靠——但在\emph{非双曲}点（有纯虚特征值，如无摩擦摆的中心）线性化会骗你（线性说“中心、周期轨”，非线性可能是稳定或不稳定焦点）。\emph{盲目不信}：以为非线性系统必须数值仿真才能知道行为——浪费了 Hartman--Grobman 提供的强大局部结论。正确姿态是：\textbf{先看平衡点是否双曲}。双曲则线性化定性可靠；非双曲则需更精细的工具（中心流形）。

\paragraph{历史。}
相平面定性分析由 Poincaré（1881--86，《微分方程定义的曲线》）开创——他放弃求闭式解，转而研究解曲线在相空间中的几何（不动点、极限环、轨迹拓扑），开创了动力系统的定性理论。Grobman（1959）与 Hartman（1960）独立证明了双曲不动点处的拓扑共轭定理，为“线性化分析”奠定严格基础\cite{khalil2002nonlinear}。

\paragraph{理论：平衡点、双曲性与二维相肖像分类。}

\begin{definition}[平衡点与双曲性]\label{def:equilibrium}
自治系统 $\dot{\mathbf{x}}=\mathbf{f}(\mathbf{x})$ 的\textbf{平衡点}（equilibrium / fixed point）$\mathbf{x}^*$ 满足 $\mathbf{f}(\mathbf{x}^*)=\mathbf{0}$（系统在此静止）。线性化矩阵 $\mathbf{A}=D\mathbf{f}(\mathbf{x}^*)$（Jacobian 在 $\mathbf{x}^*$ 处的值）的谱 $\sigma(\mathbf{A})$ 决定局部几何。$\mathbf{x}^*$ 称\textbf{双曲}（hyperbolic）当 $\sigma(\mathbf{A})\cap i\mathbb{R}=\varnothing$，即无实部为零的特征值。
\end{definition}

二维相肖像分类（$\mathbf{A}\in\mathbb{R}^{2\times2}$，迹 $\tau=\operatorname{tr}\mathbf{A}$、行列式 $\Delta=\det\mathbf{A}$）要烂熟于心：

\begin{table}[H]\centering\small
\caption{二维线性系统的不动点分类}
\begin{tabular}{llll}
\toprule
特征值类型 & 名称 & 几何行为 & 双曲? \\
\midrule
$\lambda_1,\lambda_2<0$（实，同号） & 稳定结点 (stable node) & 轨迹沿慢特征方向切入原点 & 是 \\
$\lambda_1,\lambda_2>0$（实，同号） & 不稳定结点 (unstable node) & 轨迹反向射出 & 是 \\
$\lambda_1>0>\lambda_2$（实，异号） & 鞍点 (saddle) & 稳定 / 不稳定流形成十字 & 是 \\
$\alpha\pm i\beta$，$\alpha<0$ & 稳定焦点 (stable focus) & 螺旋收敛进原点 & 是 \\
$\alpha\pm i\beta$，$\alpha>0$ & 不稳定焦点 (unstable focus) & 螺旋发散 & 是 \\
$\pm i\beta$（$\alpha=0$） & 中心 (center) & 闭合周期轨绕原点 & \textbf{否（非双曲！）} \\
\bottomrule
\end{tabular}
\end{table}

\begin{figure}[ht]
\centering
\begin{tikzpicture}[scale=1.0, >=Stealth, every node/.style={font=\footnotesize}]
  % 三个小相图：稳定焦点、鞍点、中心
  \begin{scope}[xshift=0cm]
    \draw[->,gray] (-1.3,0)--(1.3,0); \draw[->,gray] (0,-1.3)--(0,1.3);
    \foreach \a in {0,60,...,300}{
      \draw[->, main!70!black, thick, domain=\a:\a+330, samples=40, variable=\t, smooth]
        plot ({0.012*(\a+330-\t)*cos(\t)}, {0.012*(\a+330-\t)*sin(\t)});
    }
    \fill (0,0) circle (1.2pt);
    \node at (0,-1.7) {稳定焦点};
  \end{scope}
  % 鞍点：双曲线族 xy=±c，按象限分段画（避开原点奇异、范围限定在框内）
  \begin{scope}[xshift=4cm]
    \draw[->,gray] (-1.3,0)--(1.3,0); \draw[->,gray] (0,-1.3)--(0,1.3);
    \foreach \c in {0.16,0.45,0.85}{
      % 第一象限 (x>0,y>0) 与第三象限 (x<0,y<0)
      \draw[second!75!black, thick, domain=\c/1.2:1.2, samples=30, smooth, variable=\x] plot ({\x},{\c/\x});
      \draw[second!75!black, thick, domain=\c/1.2:1.2, samples=30, smooth, variable=\x] plot ({-\x},{-\c/\x});
      % 第二象限 (x<0,y>0) 与第四象限 (x>0,y<0)
      \draw[second!75!black, thick, domain=\c/1.2:1.2, samples=30, smooth, variable=\x] plot ({-\x},{\c/\x});
      \draw[second!75!black, thick, domain=\c/1.2:1.2, samples=30, smooth, variable=\x] plot ({\x},{-\c/\x});
    }
    % 稳定 / 不稳定流形（双曲线的渐近线，对角十字）
    \draw[->,thick] (-1.25,0)--(1.25,0); \draw[->,thick] (0,1.25)--(0,-1.25);
    \fill (0,0) circle (1.2pt);
    \node at (0,-1.7) {鞍点};
  \end{scope}
  % 中心
  \begin{scope}[xshift=8cm]
    \draw[->,gray] (-1.3,0)--(1.3,0); \draw[->,gray] (0,-1.3)--(0,1.3);
    \foreach \r in {0.4,0.7,1.0}{ \draw[third!75!black, thick] (0,0) circle (\r);
      \draw[->,third!75!black, thick] (\r,0.02)--(\r,-0.02);}
    \fill (0,0) circle (1.2pt);
    \node at (0,-1.7) {中心（非双曲）};
  \end{scope}
\end{tikzpicture}
\caption{三类典型不动点的相肖像：稳定焦点（螺旋收敛）、鞍点（稳定 / 不稳定流形成十字）、中心（闭合周期轨）。前两者双曲、线性化定性可靠；中心非双曲，线性化结论不可信。}
\label{fig:ode-phase-portraits}
\end{figure}

\textbf{迹--行列式判别图}：在 $(\tau,\Delta)$ 平面上，$\Delta<0$ 鞍点；$\Delta>0,\tau<0$ 稳定（$\tau^2>4\Delta$ 结点、$<4\Delta$ 焦点）；$\Delta>0,\tau>0$ 不稳定；$\tau=0,\Delta>0$ 中心（临界，非双曲）。一张图概括所有二维线性相肖像。

\begin{theorem}[Hartman--Grobman, 1959/1960]\label{thm:hartman-grobman}
若 $\mathbf{x}^*$ 是\emph{双曲}平衡点，则存在 $\mathbf{x}^*$ 的邻域与 $\mathbf{0}\in\mathbb{R}^n$ 的邻域之间的\textbf{同胚 $\mathbf{h}$}（连续双射、逆也连续），使非线性流 $\boldsymbol{\phi}_t$ 与线性流 $e^{t\mathbf{A}}$ 在邻域内\emph{局部拓扑共轭}：$\mathbf{h}\circ\boldsymbol{\phi}_t=e^{t\mathbf{A}}\circ\mathbf{h}$。
\end{theorem}
\noindent\textbf{意义}：双曲平衡点附近，非线性系统的轨迹与线性化系统的轨迹“长得一样”（可连续形变互相对应）——稳定结点还是稳定结点、鞍点还是鞍点。\emph{所以看 $\mathbf{A}$ 的特征值就能判定非线性系统的局部定性行为}。证明骨架：谱分解 $\mathbb{R}^n=E^s\oplus E^u$（稳定 / 不稳定子空间），构造共轭 $\mathbf{h}=\mathrm{id}+\mathbf{g}$，$\mathbf{g}$ 作为某函数空间上算子的不动点由 Banach 不动点定理求得；双曲性提供的 $e^{t\mathbf{A}}$ 在 $E^s$ 指数压缩、$E^u$ 指数膨胀是压缩估计的关键。

\begin{insight}[双曲性 $=$ 线性化的鲁棒性]
Hartman--Grobman 的深刻之处在于它划定了“线性化分析”的\emph{适用边界}——双曲性。双曲意味着“没有临界方向”：每个方向要么严格收缩要么严格膨胀，非线性高阶项扰动不了这个“严格”的定性结论。一旦出现纯虚特征值（中心），收缩 / 膨胀都“不严格”（中性），高阶项就能左右乾坤，线性化失效。这就是为什么控制工程师设计时偏爱“远离虚轴”的极点配置——不只是收敛快，更是定性行为稳健。注意 $\mathbf{h}$ 一般\emph{只是同胚而非微分同胚}，故 H--G 保证“定性”（拓扑）一致，但不保证“定量”（收敛速率）严格一致。
\end{insight}

\paragraph{机器人应用：倒立摆与 PD 阻尼图谱。}
\emph{(i) 倒立摆}。单摆 $\ddot\theta=-\frac{g}{\ell}\sin\theta-b\dot\theta$（$\theta$ 从竖直向下量）。下平衡点 $(0,0)$：无摩擦时特征值 $\pm i\sqrt{g/\ell}$（中心，非双曲），有摩擦 $b>0$ 时变稳定焦点（双曲）。上平衡点 $(\pi,0)$：线性化特征值 $\pm\sqrt{g/\ell}$（一正一负），\textbf{鞍点}——这就是倒立摆“天生不稳定”的数学刻画，需要主动控制。H--G 保证：上平衡点附近真实非线性摆的行为定性上就是鞍点。
\emph{(ii) PD 阻尼图谱}。线性化的 PD 闭环 $\ddot{\tilde q}+k_d\dot{\tilde q}+k_p\tilde q=0$，特征值 $\lambda=\frac{-k_d\pm\sqrt{k_d^2-4k_p}}{2}$：欠阻尼（$k_d^2<4k_p$）稳定焦点（螺旋收敛、有超调）；临界阻尼（$k_d^2=4k_p$）重根稳定结点（最快无超调）；过阻尼（$k_d^2>4k_p$）两负实根慢结点（无超调但慢）。\emph{这就是 PD 增益整定的相平面图谱}——选 $k_p,k_d$ 就是在选闭环极点位置。
\emph{(iii) 局部 LQR 的有效性}。LQR 在平衡点线性化上设计最优反馈，使闭环 $\mathbf{A}-\mathbf{B}\mathbf{K}$ 的极点都在左半平面（双曲、稳定）。H--G 保证：只要在双曲平衡点附近，这个线性最优控制器对真实非线性系统也\emph{局部}有效——这是“线性化 $+$ LQR”工程套路的理论背书。

\begin{pitfall}[概念误区：在非双曲点盲信线性化 / 把局部结论当全局]
两个相关陷阱。(a) 平衡点线性化得纯虚特征值（中心）就断言非线性也有周期轨——\emph{错}：H--G 只对双曲点成立，纯虚（非双曲）处高阶项决定真实行为（无摩擦摆加一点非线性阻尼可能变稳定焦点或不稳定焦点）。先检查双曲性，非双曲时用中心流形或直接看 Lyapunov 函数。(b) 平衡点线性化稳定就认为\emph{全局}稳定——H--G 是\emph{局部}定理（仅平衡点小邻域），全局可能有多个平衡点、极限环、吸引域边界。局部稳定用线性化，全局稳定 / 吸引域用 Lyapunov 直接法（\cref{sec:ode-lyapunov}，径向无界 $V$ 给全局结论）。
\end{pitfall}

\begin{practice}
\begin{enumerate}
  \item 写出单摆 $\ddot\theta=-\frac{g}{\ell}\sin\theta-b\dot\theta$ 的状态空间形式，求 $(0,0)$ 与 $(\pi,0)$ 处的线性化矩阵 $\mathbf{A}$，算特征值，判定类型。讨论 $b=0$ 与 $b>0$ 在 $(0,0)$ 处的区别（中心 vs 焦点）——这正是非双曲点对扰动敏感的例子。
  \item 对 PD 闭环 $\ddot{\tilde q}+k_d\dot{\tilde q}+k_p\tilde q=0$，在迹--行列式平面上定位（$\tau=-k_d$，$\Delta=k_p$），找出欠 / 临界 / 过阻尼分界线 $k_d^2=4k_p$。给定 $k_p=100$，求临界阻尼 $k_d$。说明为什么实际机器人常选略欠阻尼。
  \item 若把闭环极点配置得非常靠近虚轴（$\operatorname{Re}\lambda=-0.01$），虽然技术上双曲、稳定，实践中会有什么问题？（提示：H--G 邻域大小、对扰动 / 参数误差的鲁棒性、收敛速度。）这如何解释“极点不能只看稳定、还要有足够稳定裕度”的工程经验？
\end{enumerate}
\end{practice}

% ════════════════════════════════════════════════════════════════
\section{Lyapunov 稳定性：直接法与 LaSalle 不变原理}
\label{sec:ode-lyapunov}

\paragraph{动机：不解方程就证明“会收敛”。}
你给机械臂写了 PD$+$重力补偿，直觉上它会停在目标位形。怎么\emph{证明}它一定收敛？非线性、多自由度、耦合——闭式解算不出，线性化（\cref{sec:ode-phase}）只给局部、不给吸引域。难道只能仿真无数初值碰运气？Lyapunov 的天才在于：\textbf{构造一个“广义能量”函数 $V(\mathbf{x})$，证明它沿轨迹单调下降，就能断定系统收敛——完全不需要解方程}。这是非线性控制最强大、最通用的工具。机器人学几乎所有稳定性证明（PD、计算力矩、自适应、滑模、无源性、CBF）都是 Lyapunov 法的变体。

\paragraph{反面：没有 Lyapunov 函数会怎样。}
不会 Lyapunov 法，证明非线性系统稳定只剩两条死路：(1) 求闭式解——非线性系统几乎不可能；(2) 线性化——只给双曲平衡点的局部结论，给不出吸引域范围，也处理不了“$\dot V$ 仅半负定”的机械系统（速度负定但位置不直接负定）。许多重要系统在平衡点处非双曲或临界（保守系统、含积分器的控制），线性化彻底失效，Lyapunov 法成了唯一出路。

\paragraph{历史。}
Aleksandr Lyapunov 在 1892 年的博士论文《运动稳定性的一般问题》中创立了两套方法：\emph{第一法}（线性化，\cref{sec:ode-phase} 的源头）与\emph{第二法 / 直接法}（能量函数，本节核心）。直接法的革命性在于：不需要求解或近似方程，只需找一个合适的标量函数。LaSalle（1960）的不变原理进一步放宽了直接法的严格性要求（$\dot V$ 半负定即可），使其能直接套用于机械系统\cite{khalil2002nonlinear}。

\begin{insight}[用一个标量代理替代整个高维轨迹分析]
Lyapunov 直接法的灵魂是“用一个标量代理（能量）替代整个高维轨迹分析”。我们不追踪轨迹去哪（那需要解方程），而是追踪一个随轨迹单调变化的标量 $V$——只要 $V$ 下降且下有界，轨迹就被“挤”向 $V$ 的最小值。这与物理学“耗散系统趋向能量最小态”一致，但 Lyapunov 把它抽象成纯数学：\emph{$V$ 不必是真实物理能量，任何满足正定 $+$ 沿轨迹下降的函数都行}。这种自由度正是 Lyapunov 法威力的来源——也是难点（怎么“猜”出合适的 $V$）。
\end{insight}

\paragraph{理论：稳定性定义、直接法与 LaSalle。}
设 $\mathbf{x}^*=\mathbf{0}$ 为自治系统 $\dot{\mathbf{x}}=\mathbf{f}(\mathbf{x})$ 的平衡点（$\mathbf{f}(\mathbf{0})=\mathbf{0}$）。

\begin{definition}[三层稳定性，由弱到强]\label{def:stability}
\begin{itemize}
  \item \textbf{（Lyapunov 意义下）稳定}：$\forall\varepsilon>0,\ \exists\delta>0$，使 $\|\mathbf{x}_0\|<\delta\Rightarrow\|\mathbf{x}(t)\|<\varepsilon\ \forall t\ge0$（初值够近，轨迹一直近——“不跑远”）。
  \item \textbf{渐近稳定}：稳定\emph{且} $\exists r>0:\ \|\mathbf{x}_0\|<r\Rightarrow\mathbf{x}(t)\to\mathbf{0}$（不仅不跑远，还收敛到平衡点）。
  \item \textbf{指数稳定}：$\exists\alpha,\lambda>0:\ \|\mathbf{x}(t)\|\le\alpha\|\mathbf{x}_0\|e^{-\lambda t}$（以指数速率收敛——最强，给速率保证）。
\end{itemize}
\end{definition}
\noindent 注意“稳定”与“渐近稳定”的区别：无摩擦摆在底部是\emph{稳定但非渐近稳定}（绕着转、不跑远，但也不停下）；加摩擦后变渐近稳定（最终停下）。

\textbf{Lyapunov 直接法。} \textbf{Lyapunov 函数} $V:U\to\mathbb{R}$ 是 $C^1$、\emph{正定}（$V(\mathbf{0})=0$，$V(\mathbf{x})>0$ 对 $\mathbf{x}\neq\mathbf{0}$）的函数。其\emph{沿轨迹的导数}
\begin{equation}
  \dot V(\mathbf{x}):=\nabla V(\mathbf{x})\cdot\mathbf{f}(\mathbf{x})=\sum_i\frac{\partial V}{\partial x_i}f_i(\mathbf{x})
  \label{eq:Vdot}
\end{equation}
\emph{不需要解方程就能算}——它是 $\nabla V$ 与向量场 $\mathbf{f}$ 的内积！这是直接法“不解方程”的关键。

\begin{theorem}[Lyapunov 直接法, 1892]\label{thm:lyapunov-direct}
设 $V$ 在 $\mathbf{0}$ 的邻域上正定、$C^1$：
\begin{enumerate}
  \item 若 $\dot V(\mathbf{x})\le0$（半负定），则 $\mathbf{x}^*=\mathbf{0}$ \emph{稳定}；
  \item 若 $\dot V(\mathbf{x})<0$（$\mathbf{x}\neq\mathbf{0}$，负定），则 $\mathbf{x}^*=\mathbf{0}$ \emph{渐近稳定}；
  \item 若还\emph{径向无界}（$\|\mathbf{x}\|\to\infty\Rightarrow V(\mathbf{x})\to\infty$）且 $\dot V<0$ 全局成立，则\emph{全局渐近稳定}；
  \item 若 $\dot V\le-cV$（$c>0$），则\emph{指数稳定}（速率 $c/2$）。
\end{enumerate}
\end{theorem}
\begin{proof}[证明思路]
情形 1：$\dot V\le0\Rightarrow V(\mathbf{x}(t))\le V(\mathbf{x}_0)$，由 $V$ 正定，$V$ 的小水平集是 $\mathbf{0}$ 的邻域，轨迹困在其中，不跑远。情形 2：$\dot V<0$ 严格下降 $+$ $V\ge0$ 下有界 $\Rightarrow V(\mathbf{x}(t))\to V_\infty\ge0$；用 $\omega$-极限集论证 $V_\infty=0$（否则在 $V=V_\infty$ 处 $\dot V<0$ 矛盾），故 $\mathbf{x}(t)\to\mathbf{0}$。情形 4：$\dot V\le-cV\Rightarrow V(t)\le V(0)e^{-ct}$，$V$ 正定夹出 $\|\mathbf{x}\|$ 指数衰减。
\end{proof}

\begin{theorem}[LaSalle 不变原理, 1960]\label{thm:lasalle}
设 $\dot V\le0$（\emph{仅半负定}！）在紧正不变集 $\Omega$ 上成立。则从 $\Omega$ 出发的所有轨迹收敛到集合 $E=\{\mathbf{x}\in\Omega:\dot V(\mathbf{x})=0\}$ 中的\textbf{最大不变集 $M$}。
\end{theorem}
\noindent\textbf{意义}：直接法情形 2 要求 $\dot V<0$ 严格负定，但\emph{机械系统的 $\dot V$ 常只对速度负定}（如 $\dot V=-\dot{\mathbf{q}}^\top\mathbf{K}_d\dot{\mathbf{q}}$，在 $\dot{\mathbf{q}}=\mathbf{0}$ 但 $\mathbf{q}\neq\mathbf{q}_d$ 处为零，仅半负定）——情形 2 用不了。LaSalle 是救星：它说“轨迹最终落到 $\dot V=0$ 的最大不变集里”，再论证这个集只含平衡点，就得渐近稳定。

\begin{insight}[LaSalle 的物理图景：能量耗尽于速度为零，但只有平衡点能停住]
LaSalle 把“$\dot V<0$ 严格”换成“$\dot V\le0$ $+$ 最大不变集分析”。深层思想：轨迹最终只能停在“$V$ 不再下降”的地方（$\dot V=0$），而这些地方中真正能停住的只有\emph{不变集}；若该不变集只有平衡点，轨迹必收敛到平衡点。这恰好匹配机械系统的物理图景——能量耗散（$\dot V\le0$）只在速度为零时停止耗散（$\dot V=0$ 对应 $\dot{\mathbf{q}}=\mathbf{0}$），但速度为零、位置未到目标时系统不会“卡住”（重力 / 弹力会推它继续动，不是不变集），唯一能停住的是平衡点。
\end{insight}

\paragraph{机器人应用（极其核心）：PD$+$重力补偿。}
机器人学最著名的稳定性证明（Takegaki \& Arimoto, 1981）。控制律 $\boldsymbol{\tau}=-\mathbf{K}_p(\mathbf{q}-\mathbf{q}_d)-\mathbf{K}_d\dot{\mathbf{q}}+\mathbf{G}(\mathbf{q})$（$\mathbf{K}_p,\mathbf{K}_d\succ0$）。闭环误差动力学（令 $\tilde{\mathbf{q}}=\mathbf{q}-\mathbf{q}_d$）：$\mathbf{M}(\mathbf{q})\ddot{\tilde{\mathbf{q}}}+\mathbf{C}(\mathbf{q},\dot{\mathbf{q}})\dot{\tilde{\mathbf{q}}}+\mathbf{K}_p\tilde{\mathbf{q}}+\mathbf{K}_d\dot{\tilde{\mathbf{q}}}=\mathbf{0}$。取 Lyapunov 函数（动能 $+$ 虚拟弹性势能）
\begin{equation}
  V=\tfrac12\dot{\mathbf{q}}^\top\mathbf{M}(\mathbf{q})\dot{\mathbf{q}}+\tfrac12\tilde{\mathbf{q}}^\top\mathbf{K}_p\tilde{\mathbf{q}},
  \label{eq:pd-lyapunov}
\end{equation}
$V$ 正定（$\mathbf{M}\succ0$、$\mathbf{K}_p\succ0$）。求 $\dot V$，利用科氏矩阵的\emph{反对称性} $\dot{\mathbf{q}}^\top(\dot{\mathbf{M}}-2\mathbf{C})\dot{\mathbf{q}}=0$（即 $\tfrac12\dot{\mathbf{q}}^\top\dot{\mathbf{M}}\dot{\mathbf{q}}=\dot{\mathbf{q}}^\top\mathbf{C}\dot{\mathbf{q}}$）：
\begin{align}
  \dot V&=\dot{\mathbf{q}}^\top\mathbf{M}\ddot{\mathbf{q}}+\tfrac12\dot{\mathbf{q}}^\top\dot{\mathbf{M}}\dot{\mathbf{q}}+\tilde{\mathbf{q}}^\top\mathbf{K}_p\dot{\tilde{\mathbf{q}}}
   =\dot{\mathbf{q}}^\top\bigl(-\mathbf{C}\dot{\mathbf{q}}-\mathbf{K}_p\tilde{\mathbf{q}}-\mathbf{K}_d\dot{\mathbf{q}}\bigr)+\dot{\mathbf{q}}^\top\mathbf{C}\dot{\mathbf{q}}+\tilde{\mathbf{q}}^\top\mathbf{K}_p\dot{\mathbf{q}}\notag\\
  &=-\dot{\mathbf{q}}^\top\mathbf{K}_d\dot{\mathbf{q}}\le0\quad(\text{仅半负定！}).
  \label{eq:pd-Vdot}
\end{align}
$\dot V=0\iff\dot{\mathbf{q}}=\mathbf{0}$。在 $\dot{\mathbf{q}}=\mathbf{0}$ 上找最大不变集：$\dot{\mathbf{q}}\equiv\mathbf{0}\Rightarrow\ddot{\mathbf{q}}=\mathbf{0}$，代回闭环得 $\mathbf{K}_p\tilde{\mathbf{q}}=\mathbf{0}\Rightarrow\tilde{\mathbf{q}}=\mathbf{0}$。故最大不变集只有 $\{(\tilde{\mathbf{q}},\dot{\mathbf{q}})=(\mathbf{0},\mathbf{0})\}$。\textbf{LaSalle $\Rightarrow$ 渐近稳定，$\mathbf{q}\to\mathbf{q}_d$}。两行能量论证证明了一个非线性多自由度系统的收敛——这就是 Lyapunov $+$ LaSalle 的威力。同样的套路覆盖一大类方法：计算力矩法（反馈线性化后用线性 Lyapunov）、Slotine--Li 自适应控制（$V$ 含参数误差项）\cite{slotine1991applied}、滑模控制（$V=\tfrac12 s^2$，$\dot V\le-\eta|s|$ 给有限时间到达）、能量整形与无源性控制、控制障碍函数。

\begin{insight}[理论-工程桥接：Lyapunov 不只是事后验证，更是正向设计方法论]
为什么这些控制器都“凑得出”漂亮的 Lyapunov 函数？因为\emph{它们就是为了让某个能量函数下降而设计的}——这是逆向工程。设计流程往往是：先选 $V$（通常是动能 $+$ 某种势能 / 误差能量），再设计 $\boldsymbol{\tau}$ 使 $\dot V\le0$。PD$+$重力补偿里 $\mathbf{G}(\mathbf{q})$ 那一项的作用正是“抵消重力势能的梯度”，让 $V$ 里只剩可控的虚拟弹性势能 $\tfrac12\tilde{\mathbf{q}}^\top\mathbf{K}_p\tilde{\mathbf{q}}$。Lyapunov 是“正向设计控制器”的方法论——这是它在机器人控制中无处不在的根本原因。
\end{insight}

\begin{pitfall}[概念误区：找不到 $\dot V<0$ 就放弃，不知道用 LaSalle]
机械系统 $\dot V=-\dot{\mathbf{q}}^\top\mathbf{K}_d\dot{\mathbf{q}}$ 只半负定（$\dot{\mathbf{q}}=\mathbf{0}$ 处为零），以为证不出渐近稳定——\emph{错}。机械系统的能量耗散天然只对速度负定，$\dot V$ 几乎不可能全局严格负定，这正是 LaSalle 存在的理由。$\dot V$ 半负定时用 LaSalle：求 $\{\dot V=0\}$ 中最大不变集，证它只含平衡点，即得渐近稳定。这是机器人 Lyapunov 证明的标准套路。相关地：\emph{找不到 Lyapunov 函数 $\neq$ 系统不稳定}——定理是\emph{充分}条件，找不到可能只是没找对 $V$。
\end{pitfall}

\begin{pitfall}[概念误区：忽略径向无界 / 把稳定当渐近稳定]
两个常见混淆。(a) $\dot V<0$ 全局成立就断言全局渐近稳定，\emph{忘记检查径向无界}——若 $V$ 不径向无界（某方向 $\|\mathbf{x}\|\to\infty$ 但 $V$ 有界），水平集 $\{V\le c\}$ 可能无界、困不住轨迹，全局结论错误。声明全局渐近稳定前必须验证 $V$ 径向无界。(b) 把“稳定”等同“渐近稳定”——稳定（不跑远）比渐近稳定（收敛）弱，$\dot V\le0$ 只给稳定，$\dot V<0$ 或 LaSalle 才给渐近稳定（无摩擦摆稳定但永远振荡、不收敛）。
\end{pitfall}

\begin{practice}
\begin{enumerate}
  \item 独立复述 PD$+$重力补偿的完整 Lyapunov $+$ LaSalle 证明：写出 $V$、用反对称性算出 $\dot V=-\dot{\mathbf{q}}^\top\mathbf{K}_d\dot{\mathbf{q}}$、再用 LaSalle 论证最大不变集只含 $(\tilde{\mathbf{q}},\dot{\mathbf{q}})=\mathbf{0}$。特别说清反对称性在哪一步、为何让科氏项消失。
  \item 对标量系统 $\dot x=-x^3$，取 $V=\tfrac12 x^2$，算 $\dot V=-x^4\le0$（实际负定）。证明全局渐近稳定（验证径向无界）。再说明它\emph{不是}指数稳定（验证 $\dot V\le-cV$ 对任何 $c>0$ 在 $x\to0$ 时失败），收敛速率是多项式 $\sim1/\sqrt t$。这展示“渐近 $\neq$ 指数”。
  \item PD$+$重力补偿需精确知道 $\mathbf{G}(\mathbf{q})$。若重力模型有误差 $\hat{\mathbf{G}}=\mathbf{G}+\Delta\mathbf{G}$，平衡点会偏移（稳态误差）还是仍收敛到 $\mathbf{q}_d$？用 Lyapunov 分析新 $\dot V$ 会多出什么项。这如何解释“纯 PD 对机械臂有稳态下垂误差，而 PID 的积分项能消除它”？（提示：积分项相当于在线估计并抵消 $\Delta\mathbf{G}$。）
\end{enumerate}
\end{practice}

% ════════════════════════════════════════════════════════════════
\section{线性稳定性与 Lyapunov 方程}
\label{sec:ode-lyapeq}

\paragraph{动机：线性系统稳定性的“代数判据”与 LQR 之根。}
线性系统 $\dot{\mathbf{x}}=\mathbf{A}\mathbf{x}$ 稳定 $\iff$ $\mathbf{A}$ 的所有特征值实部为负（Hurwitz）。但有时我们不想（或不便）算特征值，而想要一个\emph{代数判据}或一个\emph{现成的} Lyapunov 函数。Lyapunov 方程 $\mathbf{A}^\top\mathbf{P}+\mathbf{P}\mathbf{A}=-\mathbf{Q}$ 同时给了这两样：它的可解性等价于 Hurwitz，它的解 $\mathbf{P}$ 直接给出二次型 Lyapunov 函数 $V=\mathbf{x}^\top\mathbf{P}\mathbf{x}$。更重要的是，\textbf{这个方程是 LQR 最优控制的数学核心}——LQR 的代数 Riccati 方程解 $\mathbf{P}$ 恰好同时是闭环系统的 Lyapunov 矩阵。理解 Lyapunov 方程，就理解了线性最优控制为何“自带稳定性证明”。

\paragraph{反面：只会算特征值会怎样。}
只靠“算特征值判稳定”有几个局限：(1) 高维系统特征值数值敏感（病态矩阵难算准）；(2) 给不出 Lyapunov 函数（无法做后续鲁棒性 / 吸引域分析）；(3) 无法自然连接到最优控制（LQR）与凸优化（LMI）。Lyapunov 方程路线绕开这些——用一个线性矩阵方程的可解性刻画稳定性，解 $\mathbf{P}$ 还能直接用于鲁棒分析与最优控制综合\cite{boyd1994linear}。

\paragraph{历史。}
Lyapunov 方程源于 Lyapunov 1892 年对线性系统的处理（第二法在线性情形的具体化）。20 世纪中叶随控制论复兴，成为 LQR（Kalman 1960）、$H_2/H_\infty$ 控制、LMI（线性矩阵不等式，Boyd 等 1994）的基石。“用矩阵不等式刻画稳定性并凸优化求解”的现代控制范式，根在这里。

\paragraph{理论：Hurwitz 矩阵与 Lyapunov 方程等价性。}

\begin{definition}[Hurwitz 矩阵]\label{def:hurwitz}
$\mathbf{A}$ 称 \textbf{Hurwitz} 当它的所有特征值有\emph{严格负实部}（$\operatorname{Re}\lambda_i<0$）。这等价于 $\dot{\mathbf{x}}=\mathbf{A}\mathbf{x}$ 渐近稳定（指数稳定）。
\end{definition}

\begin{theorem}[Lyapunov 方程等价性]\label{thm:lyapunov-eq}
以下三条等价：
\begin{enumerate}
  \item $\mathbf{A}$ 是 Hurwitz；
  \item 对\emph{任意} $\mathbf{Q}\succ0$，Lyapunov 方程 $\mathbf{A}^\top\mathbf{P}+\mathbf{P}\mathbf{A}=-\mathbf{Q}$ 有\emph{唯一}解 $\mathbf{P}$，且 $\mathbf{P}\succ0$；
  \item \emph{存在} $\mathbf{Q}\succ0$ 和 $\mathbf{P}\succ0$ 满足 $\mathbf{A}^\top\mathbf{P}+\mathbf{P}\mathbf{A}=-\mathbf{Q}$。
\end{enumerate}
\end{theorem}
\begin{proof}
\emph{(1 $\Rightarrow$ 2)} 构造解 $\mathbf{P}:=\displaystyle\int_0^\infty e^{t\mathbf{A}^\top}\mathbf{Q}\,e^{t\mathbf{A}}\,dt$。
\emph{收敛性}：$\mathbf{A}$ Hurwitz $\Rightarrow\|e^{t\mathbf{A}}\|\le Ce^{-\alpha t}$（$\alpha>0$，由谱映射 \cref{prop:exp-props}）$\Rightarrow$ 被积函数 $\le C^2\|\mathbf{Q}\|e^{-2\alpha t}$ 指数衰减，积分绝对收敛。
\emph{满足方程}：直接验证（识别被积函数是全导数）
\[
  \mathbf{A}^\top\mathbf{P}+\mathbf{P}\mathbf{A}=\int_0^\infty\frac{d}{dt}\bigl(e^{t\mathbf{A}^\top}\mathbf{Q}e^{t\mathbf{A}}\bigr)dt=\bigl[e^{t\mathbf{A}^\top}\mathbf{Q}e^{t\mathbf{A}}\bigr]_0^\infty=\mathbf{0}-\mathbf{Q}=-\mathbf{Q}
\]
（用 $\frac{d}{dt}e^{t\mathbf{A}}=\mathbf{A}e^{t\mathbf{A}}$ 与 Hurwitz 保证 $t\to\infty$ 项为 $\mathbf{0}$）。
\emph{正定性}：$\mathbf{x}^\top\mathbf{P}\mathbf{x}=\int_0^\infty(e^{t\mathbf{A}}\mathbf{x})^\top\mathbf{Q}(e^{t\mathbf{A}}\mathbf{x})\,dt>0$ 对 $\mathbf{x}\neq\mathbf{0}$（因 $e^{t\mathbf{A}}$ 可逆，$e^{t\mathbf{A}}\mathbf{x}\neq\mathbf{0}$，$\mathbf{Q}\succ0$）。故 $\mathbf{P}\succ0$。
\emph{唯一性}：Lyapunov 方程是关于 $\mathbf{P}$ 的线性方程，算子 $\mathbf{P}\mapsto\mathbf{A}^\top\mathbf{P}+\mathbf{P}\mathbf{A}$ 的特征值是 $\lambda_i+\lambda_j$，Hurwitz 下均非零，故可逆、解唯一。

\emph{(2 $\Rightarrow$ 3)} 平凡（取那个 $\mathbf{Q}$ 与对应 $\mathbf{P}$）。

\emph{(3 $\Rightarrow$ 1)} 取 Lyapunov 函数 $V=\mathbf{x}^\top\mathbf{P}\mathbf{x}$（$\mathbf{P}\succ0\Rightarrow V$ 正定）。沿轨迹
\[
  \dot V=\dot{\mathbf{x}}^\top\mathbf{P}\mathbf{x}+\mathbf{x}^\top\mathbf{P}\dot{\mathbf{x}}=\mathbf{x}^\top(\mathbf{A}^\top\mathbf{P}+\mathbf{P}\mathbf{A})\mathbf{x}=-\mathbf{x}^\top\mathbf{Q}\mathbf{x}<0\quad(\mathbf{Q}\succ0,\ \text{负定}).
\]
由 \cref{thm:lyapunov-direct} 情形 2/3，$\dot{\mathbf{x}}=\mathbf{A}\mathbf{x}$ 全局渐近稳定 $\Rightarrow$ 所有特征值实部 $<0\Rightarrow\mathbf{A}$ Hurwitz。
\end{proof}

\begin{insight}[Lyapunov 矩阵 $\mathbf{P}$ 就是“从当前状态出发到永远的累积代价”]
(1 $\Rightarrow$ 2) 构造的 $\mathbf{P}=\int_0^\infty e^{t\mathbf{A}^\top}\mathbf{Q}e^{t\mathbf{A}}dt$ 有深刻含义——它是“沿轨迹累积的二次代价”。$\mathbf{x}_0^\top\mathbf{P}\mathbf{x}_0=\int_0^\infty\mathbf{x}(t)^\top\mathbf{Q}\mathbf{x}(t)\,dt$ 正是从 $\mathbf{x}_0$ 出发、自由演化时 $\mathbf{Q}$-加权状态代价的总积分。所以 Lyapunov 矩阵 $\mathbf{P}$ 就是“从当前状态出发到永远的累积代价”——这直接通向 LQR：LQR 的 $\mathbf{P}$ 就是含控制代价的最优累积代价（值函数的 Hessian）。Lyapunov 方程与最优控制的血缘在此一目了然。
\end{insight}

\paragraph{机器人应用：LQR、观测器与 LMI。}
\emph{(i) 连续时间 LQR}。最小化 $J=\int_0^\infty(\mathbf{x}^\top\mathbf{Q}\mathbf{x}+\mathbf{u}^\top\mathbf{R}\mathbf{u})\,dt$ 受 $\dot{\mathbf{x}}=\mathbf{A}\mathbf{x}+\mathbf{B}\mathbf{u}$。最优反馈 $\mathbf{u}=-\mathbf{R}^{-1}\mathbf{B}^\top\mathbf{P}\mathbf{x}=-\mathbf{K}\mathbf{x}$，其中 $\mathbf{P}\succ0$ 满足\textbf{代数 Riccati 方程}（ARE）
\begin{equation}
  \mathbf{A}^\top\mathbf{P}+\mathbf{P}\mathbf{A}-\mathbf{P}\mathbf{B}\mathbf{R}^{-1}\mathbf{B}^\top\mathbf{P}+\mathbf{Q}=\mathbf{0}.
  \label{eq:are}
\end{equation}
关键观察：闭环 $\mathbf{A}_c=\mathbf{A}-\mathbf{B}\mathbf{R}^{-1}\mathbf{B}^\top\mathbf{P}$ 自动满足 $\mathbf{A}_c^\top\mathbf{P}+\mathbf{P}\mathbf{A}_c=-\mathbf{Q}-\mathbf{P}\mathbf{B}\mathbf{R}^{-1}\mathbf{B}^\top\mathbf{P}\prec0$。即 \textbf{Riccati 解 $\mathbf{P}$ 同时是闭环系统的 Lyapunov 矩阵}——$V=\mathbf{x}^\top\mathbf{P}\mathbf{x}$ 既是最优值函数，又是证明闭环稳定的 Lyapunov 函数。\emph{LQR 自带稳定性证明}，无需额外验证（只要 $(\mathbf{A},\mathbf{B})$ 可镇定、$(\mathbf{A},\mathbf{Q}^{1/2})$ 可检测）。
\emph{(ii) 观测器误差收敛（对偶）}。Luenberger 观测器误差 $\dot{\mathbf{e}}=(\mathbf{A}-\mathbf{L}\mathbf{C})\mathbf{e}$，设计 $\mathbf{L}$ 使 $\mathbf{A}-\mathbf{L}\mathbf{C}$ Hurwitz——与 LQR 对偶（用 $\mathbf{A}^\top,\mathbf{C}^\top$ 替 $\mathbf{A},\mathbf{B}$）。Kalman 滤波（\cref{ch:eskf}）是其最优版本（噪声协方差替代 $\mathbf{Q},\mathbf{R}$）。
\emph{(iii) LMI 与凸综合}。Lyapunov 方程的不等式版 $\mathbf{A}^\top\mathbf{P}+\mathbf{P}\mathbf{A}\prec0$、$\mathbf{P}\succ0$ 是\emph{线性矩阵不等式}（LMI），关于 $\mathbf{P}$ 凸。$H_2/H_\infty$ 控制、鲁棒稳定性都可写成 LMI 用半定规划求解——这把控制综合变成凸优化问题（详见 \cref{ch:convex}）。

\begin{insight}[理论-工程桥接：为什么实践中常调 $\mathbf{Q},\mathbf{R}$ 而非直接配极点]
LQR 通过 $\mathbf{Q},\mathbf{R}$ 间接配极点，且\emph{保证一定的稳定裕度}（连续 LQR 有 $\ge60^\circ$ 相位裕度、$[\tfrac12,\infty)$ 增益裕度的经典结果）——这是极点配置不自带的鲁棒性。$\mathbf{Q}$ 大则状态收敛快（惩罚状态偏差），$\mathbf{R}$ 大则控制省力（惩罚控制量）。调 $\mathbf{Q},\mathbf{R}$ 是在“性能 vs 控制代价 / 鲁棒性”间权衡，且 Riccati 解 $\mathbf{P}$ 自动给出 Lyapunov 证明——这种“优化即稳定”的优雅正是 LQR 风靡的原因。
\end{insight}

\begin{pitfall}[概念误区：混淆连续与离散 Lyapunov 方程 / 随意取 $\mathbf{Q}$]
两个陷阱。(a) 对离散系统 $\mathbf{x}_{k+1}=\mathbf{A}\mathbf{x}_k$ 套用连续版 $\mathbf{A}^\top\mathbf{P}+\mathbf{P}\mathbf{A}=-\mathbf{Q}$——\emph{完全错误}：离散 Lyapunov（Stein）方程是 $\mathbf{A}^\top\mathbf{P}\mathbf{A}-\mathbf{P}=-\mathbf{Q}$，对应特征值\emph{模} $<1$（Schur 稳定），而连续版对应实部 $<0$（Hurwitz）。仿真离散化后的系统务必用离散版。(b) 解 Lyapunov 方程时随便取 $\mathbf{Q}$（不定或半定）——等价性定理要求 $\mathbf{Q}\succ0$；判 Hurwitz 时取 $\mathbf{Q}=\mathbf{I}$（最简正定）解 $\mathbf{P}$ 并检查 $\mathbf{P}\succ0$。$\mathbf{Q}$ 半定（如 $\mathbf{Q}=\mathbf{C}^\top\mathbf{C}$）时另需 $(\mathbf{A},\mathbf{C})$ 可检测。
\end{pitfall}

\begin{practice}
\begin{enumerate}
  \item 独立复述 \cref{thm:lyapunov-eq} 的三向证明。重点补全 (1$\Rightarrow$2)：验证 $\mathbf{P}=\int_0^\infty e^{t\mathbf{A}^\top}\mathbf{Q}e^{t\mathbf{A}}dt$ 满足方程的那步微积分（识别被积函数是全导数 $\frac{d}{dt}(e^{t\mathbf{A}^\top}\mathbf{Q}e^{t\mathbf{A}})$）。
  \item 对 $\mathbf{A}=\bigl(\begin{smallmatrix}0&1\\-k_p&-k_d\end{smallmatrix}\bigr)$（PD 闭环状态矩阵，$k_p,k_d>0$），取 $\mathbf{Q}=\mathbf{I}$，解 $\mathbf{A}^\top\mathbf{P}+\mathbf{P}\mathbf{A}=-\mathbf{I}$ 得 $\mathbf{P}$（$2\times2$ 对称，三个未知量，解线性方程组）。验证 $\mathbf{P}\succ0$（Sylvester 判据），确认 PD 闭环 Hurwitz。
  \item LQR 闭环 $\mathbf{P}$ 同时是值函数 Hessian 与 Lyapunov 矩阵。对\emph{非线性}系统，这推广为：最优值函数 $V^*(\mathbf{x})$（满足 Hamilton--Jacobi--Bellman 方程）是一个 Lyapunov 函数。论证为什么“最优值函数沿最优轨迹下降”自动给出闭环稳定性？这如何启发“用神经网络近似值函数”的强化学习（actor-critic）也隐含稳定性结构？
\end{enumerate}
\end{practice}

% ════════════════════════════════════════════════════════════════
\section{从平衡点到吸引子：定性图景}
\label{sec:ode-attractor}

\paragraph{动机：机器人的“稳态”不止平衡点。}
前几节的稳定性都围绕\emph{平衡点}（机器人静止站住）。但许多机器人行为的“稳态”不是静止：双足步态是一条\emph{周期轨道}（极限环），CPG（中枢模式发生器）振荡器输出周期信号，甚至某些自由漂浮系统长期演化到一张不变环面。要刻画这些“系统长期被吸引到哪”的对象，需要比“平衡点”更一般的概念——\textbf{吸引子}（attractor）。本节给出极限集与吸引子的定性图景，作为后续步态分析的入口（不展开 Poincaré 映射 / Floquet 的完整理论，那属于专门的周期轨道分析）。

\paragraph{理论：$\omega$-极限集、吸引子与极限环。}

\begin{definition}[$\omega$-极限集与吸引子]\label{def:attractor}
轨迹 $\boldsymbol{\phi}_t(\mathbf{x}_0)$ 的 \textbf{$\omega$-极限集} $\omega(\mathbf{x}_0)$ 是它“最终反复趋近”的点集：$\mathbf{p}\in\omega(\mathbf{x}_0)$ 当存在 $t_k\to\infty$ 使 $\boldsymbol{\phi}_{t_k}(\mathbf{x}_0)\to\mathbf{p}$。一个紧不变集 $\mathcal{A}$ 称\textbf{吸引子}，若存在其邻域 $U$（\emph{吸引域}）使从 $U$ 出发的所有轨迹的 $\omega$-极限集都落在 $\mathcal{A}$ 内，且 $\mathcal{A}$ 不含更小的此种子集（极小性）。
\end{definition}
\noindent 最简单的吸引子是\emph{渐近稳定平衡点}（吸引子是单点）。下一层是\textbf{极限环}（limit cycle）：一条\emph{孤立}的闭轨道，邻近轨迹螺旋地卷向（稳定极限环）或卷离（不稳定）它——这正是双足周期步态、振荡器的数学对象。再往上有不变环面（拟周期）乃至奇怪吸引子（混沌），但二维自治系统\emph{不可能混沌}（Poincaré--Bendixson 定理：有界轨迹的 $\omega$-极限集若不含平衡点，必是闭轨道）——这是 CPG 设计偏爱二维振荡器的深层原因。

\begin{figure}[ht]
\centering
\begin{tikzpicture}[scale=1.0, >=Stealth, every node/.style={font=\footnotesize}]
  \draw[->,gray] (-2,0)--(2,0) node[right]{$x_1$};
  \draw[->,gray] (0,-2)--(0,2) node[above]{$x_2$};
  % 极限环（粗椭圆）
  \draw[very thick, main!75!black] (0,0) ellipse (1.3 and 1.0);
  \node[main!70!black] at (1.55,1.15) {极限环 $\Gamma$};
  % 内部螺旋向外卷向极限环
  \draw[->, second!75!black, thick, domain=0:520, samples=80, variable=\t, smooth]
    plot ({(0.18+0.0019*\t)*cos(\t)}, {(0.14+0.0015*\t)*sin(\t)});
  % 外部螺旋向内卷向极限环
  \draw[->, third!75!black, thick, domain=0:360, samples=70, variable=\t, smooth]
    plot ({(1.85-0.0013*\t)*cos(\t)}, {(1.5-0.0011*\t)*sin(\t)});
  \fill (0,0) circle (1.0pt) node[below left]{\scriptsize 不稳定焦点};
\end{tikzpicture}
\caption{稳定极限环 $\Gamma$：内部轨迹（从不稳定焦点附近）向外螺旋、外部轨迹向内螺旋，都卷向同一条孤立闭轨道。极限环是双足周期步态、CPG 振荡器的数学原型——它的“稳定性”指邻近轨迹是否被吸引，需用降一维的 Poincaré 映射刻画。}
\label{fig:ode-limit-cycle}
\end{figure}

\begin{insight}[极限环的稳定性是“降一维”的稳定性]
平衡点的稳定性看 Jacobian 特征值（\cref{sec:ode-phase}）；极限环的稳定性\emph{不能}这样看（沿轨道方向的扰动既不增也不减——总有一个零特征指数）。正确做法是取一张横截于轨道的截面（Poincaré 截面），看轨迹每绕一圈打在截面上的落点序列 $\mathbf{x}_{k+1}=\mathbf{P}(\mathbf{x}_k)$（Poincaré 映射）是否收敛到不动点——这把“连续时间极限环稳定性”降成“离散映射不动点稳定性”，维数减一。双足步态分析（混合零动力学 HZD）正是这套思想：步态是 Poincaré 映射的稳定不动点。
\end{insight}

\begin{pitfall}[概念误区：用平衡点工具分析极限环]
看到机器人“稳定地走”就想找平衡点、算 Jacobian 特征值——\emph{方法错位}。周期步态没有静止平衡点（机器人一直在动），它的“稳态”是极限环。在轨道上取一点算 Jacobian，会得到一个零特征值（沿轨道切向，对应“在轨道上前后挪一点仍在轨道上”），其余特征值（Floquet 乘子 / Poincaré 映射的特征值）才刻画横向稳定性。分析极限环必须用 Poincaré 映射 / Floquet 理论，而非平衡点线性化。
\end{pitfall}

\begin{practice}
\begin{enumerate}
  \item 范德波尔振荡器 $\ddot x-\mu(1-x^2)\dot x+x=0$（$\mu>0$）有唯一稳定极限环。写出其状态空间形式，求平衡点 $(0,0)$ 处的线性化并判断它是不稳定焦点（$0<\mu<2$）。解释为什么“原点不稳定 $+$ 大幅振荡被阻尼” 共同迫使轨迹卷向一条中等幅度的极限环。
  \item 为什么二维自治系统不可能有奇怪吸引子（混沌）？用 Poincaré--Bendixson 定理的陈述论证：有界轨迹要么趋于平衡点、要么趋于闭轨道。这对“用二维 CPG 生成稳定周期步态信号”意味着什么保证？
  \item 一条稳定极限环的 Poincaré 映射不动点 $\mathbf{x}^*$ 满足 $\|D\mathbf{P}(\mathbf{x}^*)\|$ 的所有特征值（Floquet 乘子，不含那个平凡的沿轨道方向）模 $<1$。把这一判据与 \cref{def:hurwitz} 的 Hurwitz（连续）、上面陷阱提到的 Schur（离散，模 $<1$）对应起来，说明“极限环稳定 $\leftrightarrow$ Poincaré 映射 Schur 稳定”。
\end{enumerate}
\end{practice}

% ════════════════════════════════════════════════════════════════
\section*{本章常见误解汇总}
\begin{table}[H]\centering\small
\begin{tabular}{p{0.46\textwidth} p{0.46\textwidth}}
\toprule
\textbf{常见误解} & \textbf{正确理解} \\
\midrule
二阶系统的状态就是位置 $\mathbf{q}$ & 完整状态是 $(\mathbf{q},\dot{\mathbf{q}})$，相空间维数翻倍（\cref{eq:order-reduction}）\\
“自治”就是“不受控” & 自治指 $\mathbf{f}$ 不显含 $t$；状态反馈闭环仍自治（\cref{sec:ode-basic} 陷阱）\\
$\mathbf{f}$ 连续就保证解唯一 & 只保证存在；唯一性要 Lipschitz（\cref{ex:nonunique}）\\
Lipschitz 就是可微 / 就是连续 & 函数类 $C^1\subsetneq$ Lipschitz $\subsetneq$ 连续（Lipschitz 居中）；含 $|x|$ 处处 Lipschitz 不可微 \\
存在区间 $\delta$ 是解的真实寿命 & $\delta$ 仅是局部下界，靠延拓可更长（\cref{thm:picard}）\\
时变系统解为 $e^{\int\mathbf{A}\,ds}$ & 仅 $\mathbf{A}(t)$ 与积分交换时成立，一般须数值 $\boldsymbol{\Phi}$ \\
$e^{\mathbf{A}+\mathbf{B}}=e^{\mathbf{A}}e^{\mathbf{B}}$ 总成立 & 仅当 $\mathbf{A}\mathbf{B}=\mathbf{B}\mathbf{A}$；旋转不可交换的根源（\cref{prop:exp-props}）\\
用截断幂级数算 $e^{t\mathbf{A}}$ & 大范数数值病态；反对称用 Rodrigues 闭式（\cref{thm:rodrigues}）\\
欧拉角积分姿态没问题 & 万向锁奇异；用 $\mathbf{R}\exp([\boldsymbol{\omega}]_\times\Delta t)$（\cref{eq:attitude-integration}）\\
非双曲点也能信线性化 & H--G 只对双曲点成立，中心处高阶项决定（\cref{thm:hartman-grobman}）\\
局部稳定 $\Rightarrow$ 全局稳定 & H--G / 线性化是局部；全局须径向无界 Lyapunov \\
找不到 Lyapunov 函数 $=$ 不稳定 & 定理是充分条件；可能只是没找对 $V$ \\
$\dot V$ 半负定证不出渐近稳定 & 用 LaSalle 分析最大不变集（\cref{thm:lasalle}）\\
$\dot V<0$ 全局 $\Rightarrow$ 全局渐近稳定 & 还须 $V$ 径向无界，否则轨迹可沿某方向逃逸 \\
连续 / 离散 Lyapunov 方程通用 & 连续 $\mathbf{A}^\top\mathbf{P}+\mathbf{P}\mathbf{A}=-\mathbf{Q}$（Hurwitz）vs 离散 $\mathbf{A}^\top\mathbf{P}\mathbf{A}-\mathbf{P}=-\mathbf{Q}$（Schur）\\
用平衡点 Jacobian 分析步态 & 步态是极限环，须 Poincaré / Floquet（\cref{def:attractor}）\\
\bottomrule
\end{tabular}
\end{table}

% ════════════════════════════════════════════════════════════════
\section*{本章小结}

\begin{table}[H]\centering\small
\caption{符号表}
\begin{tabular}{lll}
\toprule
符号 & 含义 & 首次出现 \\
\midrule
$\mathbf{x}(t)\in\mathbb{R}^n$ & 状态向量（机器人中常为 $(\mathbf{q},\dot{\mathbf{q}})$） & \cref{eq:ivp} \\
$\mathbf{f}(t,\mathbf{x})$ & ODE 右端向量场 & \cref{eq:ivp} \\
$\boldsymbol{\phi}_t,\ \boldsymbol{\phi}(t,\mathbf{x}_0)$ & 相流（解作为初值与时间的函数） & \cref{eq:flow-group} \\
$\mathbf{M},\mathbf{C},\mathbf{G},\boldsymbol{\tau}$ & 惯性矩阵 / 科氏离心 / 重力 / 关节力矩 & \cref{eq:euler-lagrange} \\
$L$ & Lipschitz 常数 & \cref{eq:lipschitz} \\
$T$ & Picard 积分算子 & \cref{thm:picard} \\
$\boldsymbol{\Phi}(t),\ \boldsymbol{\Phi}(t,t_0)$ & 基本解矩阵 / 状态转移矩阵 & \cref{eq:fund-matrix} \\
$e^{t\mathbf{A}}$ & 矩阵指数（常系数基本解矩阵） & \cref{def:matrix-exp} \\
$[\boldsymbol{\omega}]_\times$ & $\boldsymbol{\omega}\in\mathbb{R}^3$ 的反对称矩阵 & \cref{thm:rodrigues} \\
$\mathbf{A}=D\mathbf{f}(\mathbf{x}^*)$ & 平衡点处的线性化（Jacobian） & \cref{def:equilibrium} \\
$\tau,\Delta$ & 二维 $\mathbf{A}$ 的迹与行列式 & \cref{sec:ode-phase} \\
$V(\mathbf{x}),\ \dot V$ & Lyapunov 函数及其沿轨迹导数 & \cref{eq:Vdot} \\
$\mathbf{P}\succ0,\ \mathbf{Q}\succ0$ & 对称正定矩阵（Lyapunov 方程） & \cref{thm:lyapunov-eq} \\
$\mathcal{A},\ \omega(\mathbf{x}_0)$ & 吸引子 / $\omega$-极限集 & \cref{def:attractor} \\
$\Gamma$ & 极限环（孤立闭轨道） & \cref{fig:ode-limit-cycle} \\
\bottomrule
\end{tabular}
\end{table}

\begin{table}[H]\centering\small
\caption{定理速查表}
\begin{tabular}{lll}
\toprule
定理 / 公式 & 一句话说明 & 对应 \\
\midrule
状态空间化 \cref{eq:order-reduction} & 高阶 $\to$ 一阶；状态 $=$ 位置 $+$ 速度 & \cref{sec:ode-basic} \\
Picard--Lindelöf \cref{thm:picard} & Lipschitz $+$ 压缩映射 $\Rightarrow$ 唯一解 & \cref{sec:ode-picard} \\
常数变易 \cref{eq:variation-constants} & 自由响应 $+$ 卷积响应 & \cref{sec:ode-linear} \\
Liouville \cref{eq:liouville} & $\operatorname{tr}\mathbf{A}$ 决定相空间体积演化 & \cref{prop:liouville} \\
矩阵指数 Jordan \cref{eq:expm-jordan} & 幂零 $\to$ 有限多项式，闭式 $e^{t\mathbf{A}}$ & \cref{prop:exp-props} \\
罗德里格斯 \cref{eq:rodrigues-ode} & $\so(3)$ 指数闭式，姿态积分内核 & \cref{thm:rodrigues} \\
Hartman--Grobman \cref{thm:hartman-grobman} & 双曲点线性化定性可靠 & \cref{sec:ode-phase} \\
Lyapunov 直接法 \cref{thm:lyapunov-direct} & $V$ 正定 $+\dot V\le0\Rightarrow$ 稳定 & \cref{sec:ode-lyapunov} \\
LaSalle 不变原理 \cref{thm:lasalle} & $\dot V$ 半负定 $+$ 最大不变集 & \cref{thm:lasalle} \\
PD$+$重力补偿 \cref{eq:pd-Vdot} & 两行 Lyapunov 证渐近稳定 & \cref{sec:ode-lyapunov} \\
Lyapunov 方程 \cref{thm:lyapunov-eq} & Hurwitz $\iff\exists\mathbf{P}\succ0$ & \cref{sec:ode-lyapeq} \\
LQR Riccati \cref{eq:are} & 解 $\mathbf{P}$ 兼值函数与 Lyapunov 矩阵 & \cref{sec:ode-lyapeq} \\
吸引子 / 极限环 \cref{def:attractor} & 稳态不止平衡点；步态 $=$ 极限环 & \cref{sec:ode-attractor} \\
\bottomrule
\end{tabular}
\end{table}

\begin{table}[H]\centering\small
\caption{知识点总表}
\begin{tabular}{clll}
\toprule
\# & 知识点 & 核心要点 & 难度 \\
\midrule
1 & ODE 与状态空间 & 一切机器人动力学都是 $\dot{\mathbf{x}}=\mathbf{f}$；状态 $=$ 位置 $+$ 速度 & \(\bigstar\bigstar\) \\
2 & 存在唯一性 & Lipschitz $+$ 压缩映射；唯一性 $=$ 仿真可复现 & \(\bigstar\bigstar\bigstar\) \\
3 & 线性系统 & 基本解矩阵、常数变易、叠加结构 & \(\bigstar\bigstar\) \\
4 & 矩阵指数 & $e^{t\mathbf{A}}$ 闭式（Jordan）；Rodrigues 接 $\SO(3)$ & \(\bigstar\bigstar\bigstar\) \\
5 & 相图 / 线性化稳定性 & 不动点分类；Hartman--Grobman & \(\bigstar\bigstar\bigstar\) \\
6 & Lyapunov 直接法 / LaSalle & 不解方程证收敛；PD 证明 & \(\bigstar\bigstar\bigstar\) \\
7 & Lyapunov 方程 & Hurwitz $\iff\mathbf{P}\succ0$；LQR 之根 & \(\bigstar\bigstar\bigstar\) \\
8 & 吸引子 / 极限环 & 稳态不止平衡点；步态 $=$ 极限环 & \(\bigstar\bigstar\bigstar\) \\
\bottomrule
\end{tabular}
\end{table}

% ════════════════════════════════════════════════════════════════
\section*{延伸阅读}
\begin{itemize}
  \item \textbf{严谨主线（非线性系统与稳定性）}：Khalil《Nonlinear Systems》\cite{khalil2002nonlinear}——存在唯一性、比较引理、Lyapunov 与 LaSalle、输入到状态稳定性，是控制视角 ODE 理论的权威教材，本章稳定性各节的完整出处。
  \item \textbf{机器人控制中的 Lyapunov 设计}：Slotine \& Li《Applied Nonlinear Control》\cite{slotine1991applied}——滑模、自适应、无源性控制的 Lyapunov 构造，PD / 自适应证明的范本。
  \item \textbf{线性系统理论}：Chen《Linear System Theory and Design》\cite{chen1999linear}——基本解矩阵、状态转移、能控 / 能观、实现理论，本章线性系统一节的展开来源。
  \item \textbf{LMI 与凸控制综合}：Boyd 等《Linear Matrix Inequalities in System and Control Theory》\cite{boyd1994linear}——把 Lyapunov 方程升级为 LMI、用半定规划求解（与 \cref{ch:convex} 衔接）。
  \item \textbf{几何 / 经典力学视角}：Arnold《Mathematical Methods of Classical Mechanics》\cite{arnold1989mathematical}——相流、辛结构、保守系统，为 \cref{ch:lie} 的流形语言与辛积分器埋下伏笔。
\end{itemize}

\section*{本章与后续章节的关系}
\begin{table}[H]\centering\small
\begin{tabular}{lll}
\toprule
后续章节 & 关系 & 本章铺垫 \\
\midrule
\cref{ch:lie} 李群与李代数 & $\so(3)\to\SO(3)$ 指数映射 & 矩阵指数、罗德里格斯（\cref{sec:ode-expm}）\\
\cref{ch:lyapunov-basics} Lyapunov 基础 & 稳定性定义与直接法 & \cref{thm:lyapunov-direct}、LaSalle \\
\cref{ch:lyapunov-stability} Lyapunov 稳定性 & 机器人控制律的稳定性证明 & PD 证明、Lyapunov 方程（\cref{sec:ode-lyapeq}）\\
\cref{ch:eskf} 误差状态滤波 & 连续时间状态 / 协方差传播 & 状态转移矩阵、$e^{t\mathbf{A}}$（\cref{sec:ode-linear}）\\
\cref{ch:linalg} 线性代数 & 特征值 / Jordan / 正定二次型 & 被本章 $e^{t\mathbf{A}}$、Lyapunov 反复调用 \\
\cref{ch:convex} 凸优化 & LMI / 半定规划 & Lyapunov 不等式 $\mathbf{A}^\top\mathbf{P}+\mathbf{P}\mathbf{A}\prec0$ \\
\bottomrule
\end{tabular}
\end{table}

\section*{故障排查手册}
\begin{enumerate}
  \item \textbf{症状}：仿真换求解器 / 步长\emph{轨迹大变}。\textbf{排查}：检查 $\mathbf{f}$ 是否含 $\operatorname{sign}$ / 硬接触等不连续项（破坏 Lipschitz、唯一性，\cref{ex:nonunique}）；改 Filippov / 软接触从建模层恢复唯一性，而非 debug 代码。
  \item \textbf{症状}：积分器\emph{发散到无穷}。\textbf{原因}：可能是有限时间爆破（如 $\dot x=x^2$ 型增长，\cref{sec:ode-picard}）而非 bug。\textbf{对策}：检查 $\mathbf{f}$ 增长是否超线性；加状态饱和 / 截断。
  \item \textbf{症状}：姿态积分\emph{在某姿态跳变 / 发散}。\textbf{原因}：欧拉角万向锁。\textbf{对策}：改用 $\mathbf{R}\leftarrow\mathbf{R}\exp([\boldsymbol{\omega}]_\times\Delta t)$（\cref{eq:attitude-integration}）或四元数积分。
  \item \textbf{症状}：朴素积分多步后 $\mathbf{R}$\emph{不再正交}（$\det\neq1$）。\textbf{原因}：一阶 Euler 不保流形。\textbf{对策}：用指数映射积分（内蕴），或周期 SVD / Gram--Schmidt 重正交化。
  \item \textbf{症状}：解 Lyapunov 方程得到的 $\mathbf{P}$\emph{非正定}。\textbf{排查}：确认用了 $\mathbf{Q}\succ0$ 且系统对应正确（连续用 $\mathbf{A}^\top\mathbf{P}+\mathbf{P}\mathbf{A}$、离散用 $\mathbf{A}^\top\mathbf{P}\mathbf{A}-\mathbf{P}$）；若 $\mathbf{A}$ 非 Hurwitz 则本无正定解（\cref{thm:lyapunov-eq}）。
  \item \textbf{症状}：用平衡点 Jacobian 分析周期步态\emph{得到一个零特征值、结论无意义}。\textbf{原因}：步态是极限环不是平衡点。\textbf{对策}：改用 Poincaré 映射 / Floquet 乘子（\cref{sec:ode-attractor}）。
\end{enumerate}
